\chapter{Categorical E-Graphs for \texorpdfstring{$\lambda$}{lambda}-Calculi}
\label{chapter:lambda}

As noted in~\cref{chapter:introduction}, traditional e-graphs encounter issues in the context of $\lambda$-calculi, mainly because the former were not designed with variable binding and substitution in mind.
It is well-known that different variations of $\lambda$-calculi are modelled by different flavours of closed monoidal categories.
The morphisms in the latter can be represented by string diagrams of a particular kind, as used by Ghica and Zanasi~\cite{ghica2024stringdiagramslambdacalculifunctional}.
We recapitulate the key properties of such string diagrams with respect to variables and binding and argue that this representation may be useful for performing equality saturation on the $\lambda$-calculus with explicit substitution terms by showing that particular equalities of such calculus are absorbed by this representation. 
We then develop a combinatorial representation of the version of these diagrams extended with $\catname{SLat}$-structure, and prove the equivalence between rewriting of closed-monoidal terms and rewriting of these combinatorial objects.

\section{String diagrams for closed monoidal categories}
\begin{definition}%
	\label{def:closed}
	A (right) closed monoidal category is a monoidal category $\catname{C}$ such that for
	every pair of objects $X,A$ there is an object $X \multimap A$ and a morphism $\ev_{X,A} : (X \multimap A) \otimes X \to
		A$, and for every triple of objects $A,X,B$ there is an operation $\Lambda_{A,X,B} : \catname{C}(A \otimes X,B) \to
		\catname{C}(A,X \multimap{} B)$ such that for all $f : A \otimes X \to B$ and $g : Z \to A$ the following equations hold:
	\begin{equation}
		f = (\Lambda_{A,X,B}(f) \otimes \id_{X}) \fatsemi \ev_{X,B},
  \label{eq:lambda_1}
  \tag{1}
  \end{equation}
  \begin{equation}
		\id_{X \multimap{} B} = \Lambda_{X \multimap{} B,X,B}(\ev_{X,B}),
  \label{eq:lambda_2}
  \tag{2}
  \end{equation}
  \begin{equation}
		\Lambda_{Z,X,B}((g \otimes \id_{X}) \fatsemi f) = g \fatsemi \Lambda_{A,X,B}(f).
  \label{eq:lambda_3}
  \tag{3}
  \end{equation}
\end{definition}

This is equivalent to defining a (right) closed monoidal category as one admitting an adjunction $- \otimes X \dashv X \multimap{} -$ for every object $X$.
If the tensor is symmetric, left and right closedness coincide, and we simply call the category \emph{closed}.
Note the asymmetry: $\ev_{X,B} : (X \multimap B) \otimes X \to B$ is the counit of the adjunction, whereas $\Lambda_{A,X,B}$ is the \emph{currying} isomorphism of the adjunction given by $\Lambda_{A,X,B}(f) = \eta_{A} \fatsemi (X \multimap f)$.

In addition to monoidal closed categories being models, which is useful for model-theoretic study of $\lambda$-calculi, we can pick up on some string diagrammatic wisdom to get better representations of $\lambda$-terms.
String diagrams for closed monoidal categories are given by the combined syntax in~\cref{fig:lambda-string-combined}---the syntax used by Ghica and Zanasi in~\cite{ghica2024stringdiagramslambdacalculifunctional}.
Note that wires are labelled with types.

\begin{figure}
% \begin{subfigure}{0.6\textwidth}
%     \[
%     \adjustbox{max width=\textwidth}{
%     \tikzfig{./figures/introduction/string-diagrams-syntax}
%     }
%     \]
%     \subcaption{String diagrams syntax for symmetric monoidal categories.}
% \end{subfigure}
% \hfill
% \begin{subfigure}{0.35\textwidth}
%   \[
%   \adjustbox{max width=\textwidth}{
%     \tikzfig{./figures/introduction/lambda-string}
%   }
%   \]
%     \subcaption{Extension of string diagrammatic syntax with $\lambda$-abstraction ($\Lambda$) and application $\ev$.}
% \end{subfigure}
\[
\adjustbox{max width=\textwidth}{
  \tikzfig{./figures/lambda/lambda-strings}
}
\]
\caption{String diagrams for closed monoidal categories; dashes stand for arbitrary number of wires and appropriate labels.}
\label{fig:lambda-string-combined}
\end{figure}

The equations \eqref{eq:lambda_1}--\eqref{eq:lambda_3} are represented as the equations between string diagrams in~\cref{fig:lambda-string-eq},
\begin{figure}
  \[
\adjustbox{max width=\textwidth}{
\tikzfig{./figures/lambda/lambda_eq_1},
}
\]
\caption{String diagram equations for the closed monoidal category axioms.}
\label{fig:lambda-string-eq}
\end{figure}
where the wires are labelled by the types of the corresponding objects and the lower boundary of the diagram represents the context of the represented $\lambda$-term.
If one wishes to further model $\lambda$-calculi where the context admits contraction and weakening, the string diagrams are extended with natural copy $\comonoid$ and $\counit$ delete generators.

The first representational advantage of string diagrams is that they are a \emph{nameless} representation, that is, $\alpha$-equivalent terms are represented by the same string diagram.
For example, the string diagram below represents the whole $\alpha$-equivalence class of $\lambda$-terms, including $\lambda f . f((\lambda x. f x) 2)$ and $\lambda g . g((\lambda y . g y) 2)$.
\[
    \tikzfig{./figures/introduction/closed-string-diagram-example}
\]
It is worth noticing that the same class of terms may be represented by a single $\lambda$-term using de Bruijn indices, namely the term $\lambda . \%0((\lambda . \%1 \%0) 2)$, however the different occurrences of the binder corresponding to $f$ use different indices, which might not be desirable.

Another advantage is the conceptual simplicity of $\beta$-reduction, the rule for which is the first equation in~\cref{fig:lambda-string-eq}.
In order to $\beta$-reduce a string diagrammatic application, it is sufficient to re-attach the matching wires.
This principle of using wires for the variables goes further.
Traditionally, substitution of $\lambda$-terms is made a first-class term former with its own set of equations.
Different sets of equations result in different $\lambda$-calculi.
Below we consider the \emph{linear substitution calculus} $\lambda_{\text{lsub}}$ of Accattoli et al.~\cite[Sec.~2]{accattoli2014nonstandard}.

We summarise the linear substitution calculus as the $\lambda$-calculus with an additional term former representing explicit substitution; i.e.\ terms are generated by the grammar
\begin{equation}
	t,u \Coloneqq x \mid t u \mid \lambda x. t \mid t [x / u],
  \label{eq:lambda_lsub_grammar}
\end{equation}
where $x$ is a variable drawn from some countably infinite set, $t u$ is an application, and $\lambda x. t$ is an abstraction, as in standard $\lambda$-calculus.
The term $t [x / u]$ is an explicit substitution of the term $u$ for the variable $x$ in the term $t$.

The notions of context $C[-]$, bound and free variable are as usual, with the addition that $\textsf{fv} (t [x/u]) \defeq \textsf{fv} (t) \setminus \{ x \} \cup \textsf{fv} (u)$.
Write $C \llbracket t \rrbracket$ for the term given by substitution of the hole in the context $C[-]$ for the term $t$, under the additional restriction that the free variables of $t$ are not captured (via abstraction or explicit substitution) by $C$.
In addition to the $\beta$- and $\eta$-reduction rules of standard $\lambda$-calculus, there are reduction rules to account for this explicit substitution which \enquote{perform} the substitution or commute it with $\beta$-reduction:
\begin{equation}
\label{eq:lambda_lsub_reduction}
\begin{aligned}
	(\lambda x. t) L u              \quad & \to_{d\beta} \quad t[x/u] L,                                                \\
	C \llbracket x \rrbracket [x/u] \quad & \to_{ls} \quad C \llbracket u \rrbracket [x/u],                             \\
	t [x/u]                         \quad & \to_{gc} \quad t,                               & x \notin \textsf{fv} (t),
\end{aligned}
\end{equation}
where $L$ is a finite list of substitutions of the form $[x/u] [y/v] \ldots$.
In particular, note that it is $ls$ and $gc$ reduction rules that perform the actual computation computing a substitution-free term.

In programming terms, this explicit substitution can be used to model \textbf{let}-bindings and sharing.
For instance, the ML-style term $\textbf{let} \; x = 1 \; \textbf{in} \; \text{plus} (x, x)$ is formally represented as the term $\text{plus} (x, x) [x / 1]$ in this calculus.
In this case, the variable $x$ is valued in the constant $1$, but one could imagine instead that $x$ is to be substituted for the result of some expensive computation which we do not desire to repeat twice unnecessarily---in this sense, the sharing of both operands of the operation $\text{plus}$ is captured by the explicit substitution.

\emph{Graphical equivalence} $\sim$ of $\lambda_{\text{lsub}}$-terms is augmentation of $\alpha$-equi\-va\-lence of terms with the additional stipulations that
% \begin{equation}
% \begin{aligned}
\begin{align}
	t [x/u] [y/v] \quad        & \sim \quad t [y/v] [x/u],        & x \notin \textsf{fv} (v) \land y \notin \textsf{fv} (u),
  \tag{$1$}
  \label{eq:non_interacting_subst}\\
	(\lambda y. t) [x/u] \quad & \sim \quad \lambda y. (t [x/u]), & y \notin \textsf{fv} (u),                             
  \tag{$2$}
  \label{eq:subst-lambda}\\
	(t v) [x/u] \quad          & \sim \quad t [x/u] v,            & x \notin \textsf{fv} (v),
  \tag{$3$}
  \label{eq:subst-application}
% \end{aligned}
\end{align}
% \label{eq:gs_equivalence}
% \tag{$\ast$}
% \end{equation}
closed under contexts, transitivity, symmetry, and reflexivity.
Put simply, it is the stipulation that \textbf{let}-bindings may be freely reordered without changing the meaning of the term.
The $\lambda^\sim_{\text{lsub}}$-calculus is given by these graphical equivalence classes of terms with the reduction relation defined by that of the $\lambda_{\text{lsub}}$-calculus modulo graphical equivalence.
Importantly, equivalence~\eqref{eq:subst-lambda} is different from $ls$ rule: first of all, it is not oriented, but more substantially, it does not perform any computation---it merely moves the substitution node inside abstraction.

This calculus is explicitly designed to capture the representation of terms as proof nets: for two terms $t, u$ in the $\lambda_{\text{lsub}}$-calculus, $t \sim u$ if and only if both map to the same proof net~\cite{accattoliJumpingBoxGraphical2011}.
% By further extending the equivalence relation $\sim$, taking the symmetric closure by $\to_{ls}$ and $\to_{gc}$ and thus identifying terms under elimination of explicit substitution, we determine a \emph{string diagram equivalence} $\approx$ of $\lambda_{\text{lsub}}$-terms, determining a $\lambda^\approx_{\text{lsub}}$-calculus.
\begin{proposition}[{\cite[Sec.~5.1.1]{ghica2024stringdiagramslambdacalculifunctional}}]
\label{prop:lambda_string_diagram_equivalence}
	Let \emph{string diagram equivalence} $\approx$ of $\lambda_{\text{lsub}}$-terms be the extension of graph equivalence $\sim$ by taking the symmetric closure by $\to_{ls}$ and $\to_{gc}$, thus identifying terms under elimination of explicit substitution.
	It defines the corresponding $\lambda^\approx_{\text{lsub}}$-calculus.
	We then have that $t \approx u$ if and only if both map to the same string diagram.
\end{proposition}
For example, both sides of the equivalence~\eqref{eq:non_interacting_subst} are represented by the same string diagram on the left and both sides of the equivalence~\eqref{eq:subst-application} are represented by the same string diagram on the right:
\begin{equation}
\tikzfig{./figures/lambda/gs_equivalence_example_1} \qquad \qquad \tikzfig{./figures/lambda/gs_equivalence_example_2}~.
\label{eq:string_diagram_equivalence_example}
\tag{$\ast$}
\end{equation}
% The second equivalence is the diagram on the right in the third equation in~\cref{fig:lambda-string-eq}: recall that substitution is a wire, so the wire that goes inside the abstraction represents exactly this equivalence.
% The computation then happens by moving the node $g$ in this equation inside the abstraction.
The equivalence~\eqref{eq:subst-lambda} is the third equation in~\cref{fig:lambda-string-eq}.

The string diagram equivalence $\approx$ as defined in~\cref{prop:lambda_string_diagram_equivalence} concerns string diagrams as in~\cref{fig:lambda-string}, equipped with natural copy $\comonoid$ and delete $\counit$ and quotiented by their respective equations plus the third equation in~\cref{fig:lambda-string-eq}.
Thus, this notion of equality of string diagrams is coarser than them just being the same graph.
However, the equivalences~\eqref{eq:non_interacting_subst} and~\eqref{eq:subst-application} are indeed rendered as equal graphs.
We illustrate the practical implications of performing equality saturation under these equivalences next.

% This allows us to exploit the string diagram setting of e-hypergraphs to absorb many isomorphisms of symmetric monoidal closed $\catname{SLat}$-categories (e.g.\ associators, unitors, etc.) which generate the \enquote{isotopy of string diagrams}; even in the graphical equivalence relation $\sim$, while~\cref{eq:subst-lambda} is validated in every symmetric monoidal closed $\catname{SLat}$-category, it is not so by isotopy, but rather the adjunction $ - \otimes A \dashv A \multimap -$. In other words, terms floating freely through the boundaries of rounded boxes is validated by our semantics, but produces concretely different e-hypergraphs --- this is the purpose of rectification by quotienting along $\mathcal{S}$ in~\cref{prop:quotient-structural}, and manifests as additional DPO rewrites to be imposed.} of the reductions arising due to the equivalence $\approx$ into our representation rather than explicitly providing them as equations attached to a closed symmetric monoidal theory presenting the $\lambda^\approx_{\text{lsub}}$-calculus.

% Before we proceed to demonstrate this, we first wish to discuss the relevance of this calculus in the domain of e-graphs.
% E-graphs excel as a toolkit for compiler optimisation, facilitating the rewriting of programs into more efficient forms by equality saturation.
% The theoretical study of this is the field of term rewriting, where programs are often simplistically represented as (the abstract syntax of) terms in the $\lambda$-calculus.
% However, in a modern optimising compiler (e.g.\ LLVM~\cite{lattnerLLVMCompilationFramework2004}, GNU Compiler Collection~\cite{merrillGenericGimpleNew2003}), programs are represented as some sort of graph-based intermediate representation (IR) rather than a syntax tree.
% Here, basic blocks of instructions\footnote{Maximal sequences of instructions which are unconditionally executed together determined by program control flow.} form the nodes, and terminator instructions (e.g.\ jumps, branches, and other control flow operations) form the edges.
% This allows for optimisations and analyses based on graphical structure, e.g.\ reachability analysis to perform dead-code elimination, or data-flow analysis to perform constant folding.
% In essence, this is much closer to $\lambda^\approx_{\text{lsub}}$-calculus: the string diagram equivalence there is abstracting away from the particular syntax of a term, in much the same way that LLVM IR abstracts away from the raw syntax of the source program.
% Thus, an understanding of e-graphs in this context is valuable for the applicability of equality saturation techniques with respect to graph-based IRs which dominate the landscape of modern compiler design.

% We will define a monoidal signature whose closed symmetric monoidal theory is the $\lambda^\approx_{\text{lsub}}$-calculus extended by tupling of terms and flat natural numbers equipped with basic arithmetic operations.
% Let $\Sigma \coloneq (\{ \mathbf{N} \}, \{ \comonoid, \counit \} \cup \mathbb{N} \cup \{ \text{plus}, \text{mult} \}, t)$ where
% \begin{align*}
% 	t\colon \{ \comonoid, \counit \} \cup \mathbb{N} \cup \{ \text{plus}, \text{mult} \} & \to {\{ \mathbf{N} \}}^* \times {\{ \mathbf{N} \}}^*                                                    \\
% 	t(n)                                                                                 & = ((), (\mathbf{N})),                                & n \in \mathbb{N},                                \\
% 	t(o)                                                                                 & = ((\mathbf{N}, \mathbf{N}), (\mathbf{N})),          & o \in \{ \text{plus}, \text{mult} \}, \\
% 	t(\comonoid)                                                                         & = ((\mathbf{N}), (\mathbf{N}, \mathbf{N})),                                                               \\
% 	t(\counit)                                                                           & = ((\mathbf{N}), ()).
% \end{align*}
% The operations that the signature $\Sigma$ avails are copy and delete, denoted by $\comonoid$ and $\counit$ respectively, constants for each natural number, and the binary arithmetic operations $+$ and $\times$.
% We interpret the $\Sigma$-terms $\llbracket \tau \rrbracket$ as follows: operations as themselves, $\id_{I}$ and $\id_{\mathbf{N}}$ as the identity function $\lambda x. x$ on the empty tuple $()$ and $\mathbf{N}$ respectively, $\sym_{\mathbf{N}, \mathbf{N}}$ as the swap function $\lambda (x, y). (y, x)$, composition $\tau_1 \fatsemi \tau_2$ as the composition function $\lambda x. \llbracket \tau_2 \rrbracket \left(\llbracket \tau_1 \rrbracket x\right)$, and $\tau_1 \otimes \tau_2$ as tupling $\left(\llbracket \tau_1 \rrbracket, \llbracket \tau_2 \rrbracket\right)$.

% Note that, in contrast to conventional presentations of the $\lambda$-calculus as a language of (slotted~\cite[Example 1]{slotted-egraphs}) e-graphs~\cite[Figure 10]{EggPaper}, here there is no encoding of the term formers for application, abstraction, and explicit substitution.
% Instead, these constructs are handled ambiently by our semantic setting: any closed symmetric monoidal category (and in particular, $\catname{CS}(\Sigma, \mathcal{E})$ for any set of equations $\mathcal{E}$) admits an internal language in the multiplicative fragment of intuitionistic linear logic (a particular type of linear $\lambda$-calculus)~\cite[§ 1.7.3]{abramskyIntroductionCategoriesCategorical2010}.
% In other words, with respect to~\cref{fig:egraph-strings}, this allows us to \emph{natively} represent abstraction by the rounded boxes of ($\Lambda(f)$), and application by the $\textsf{ev}$ operation with $\beta$- and $\eta$-conversion arising via equality of morphisms in the category (that is, without having to provide them as data in the form of equations in $\mathcal{E}$).
% For explicit substitution, we note that the linearity is not desirable; this is why the copy-delete comonoid $(\comonoid, \counit)$ is present: by Fox's theorem~\cite{foxCoalgebrasCartesianCategories1976}, a Cartesian category is precisely a monoidal one admitting such a \emph{natural} copy-delete structure.
% Thus, we must endow $\mathcal{E}$ with equations stating that $(\comonoid, \counit)$ forms a comonoid, and moreover that it is natural: $\comonoid$ behaves uniformly as a copy map, and $\counit$ behaves as deletion.

% \begin{figure}
% 	\centering
% 	\begin{subfigure}{0.3\linewidth}
% 		\begin{adjustbox}{scale=0.8}
% 			\begin{lstlisting}[language=ML]
% let x = 1 in
% let y = 2 in
% plus(x, y)
%                 \end{lstlisting}
% 		\end{adjustbox}
% 		\hrule
% 		\begin{adjustbox}{scale=0.8}
% 			\begin{lstlisting}[language=ML]
% let y = 2 in
% let x = 1 in
% plus(x, y)
%                 \end{lstlisting}
% 		\end{adjustbox}
% 		\hrule
% 		\begin{adjustbox}{scale=0.8}
% 			\begin{lstlisting}[language=ML]
% plus(1, 2)
%                 \end{lstlisting}
% 		\end{adjustbox}
% 		\caption{Three programs for $1 + 2$}%
% 		\label{fig:plus-1-2-programs}
% 	\end{subfigure}
% 	\hspace{1em}
% 	\begin{subfigure}{0.2\linewidth}
% 		\[
% 			\begin{adjustbox}{scale=0.8}$
% 					\begin{aligned}
% 						        & \text{plus}(x, y)[y/2][x/1] \\
% 						\approx &                             \\
% 						        & \text{plus}(x, y)[x/1][y/2] \\
% 						\approx &                             \\
% 						        & \text{plus}(1, 2)
% 					\end{aligned}
% 				$\end{adjustbox}
% 		\]
% 		\caption{As terms in $\lambda_{\text{lsub}}$}%
% 		\label{fig:plus-1-2-terms}
% 	\end{subfigure}
% 	\hspace{2em}
% 	\begin{subfigure}{0.3\linewidth}
% 		\centering
% 		\scalebox{0.8}{
% 			\tikzfig{./figures/lambda/plus-1-2}
% 		}
% 		\caption{String-diagrammatic representation}%
% 		\label{fig:plus-1-2-string-diagram}
% 	\end{subfigure}
% 	\caption{$\approx$-equivalent terms}%
% 	\label{fig:plus-1-2}
% \end{figure}

% \cref{fig:plus-1-2-programs} shows some programs which are semantically equivalent but syntactically distinct.
% As alluded to in~\cref{sec:vs-e-graphs-with-binding}, the traditional (slotted) e-graph approach to modelling the equivalence of these terms (\cref{fig:plus-1-2-terms}) relies on explicitly encoding substitution operations as e-graph nodes, and their behaviour (the string diagram equivalence relation $\approx$) as rewrite rules, as in~\cref{fig:e-graph-substitution} --- not only is this more complicated, but it also leads to a blow-up in the state space required to perform equality saturation in any extension of the $\lambda^\approx_{\text{lsub}}$-calculus, because we are really rewriting $\lambda_{\text{lsub}}$-terms and computational resources are spent rewriting between equivalent terms (an empirical analysis is provided in~\cref{sec:e-graph-explicit-subst}).
% In contrast, in our framework, $\lambda^\approx_{\text{lsub}}$-terms are directly modelled as string diagrams (\cref{fig:plus-1-2-string-diagram}\footnote{Note that dashed boxes are elided for convenience here because no rewriting has occurred.}).
% We do however note that the slotted approach is more flexible, in the sense that it may model more theories with binders (e.g.\ $\pi$-calculus), whereas our approach is essentially specialised to $\lambda$-calculus, relying heavily on the metatheoretic internal language of the ambient setting.
\section{Equality saturation for $\lambda$-calculi}

The above concepts of $\alpha$-equivalence and substitution are aspects that were notoriously difficult to support in traditional e-graphs.
The usual solution was to encode $\lambda$-terms using de Bruijn indices, which we noted is not ideal because it can obscure sharing.
Moreover, performing $\beta$-reduction on de Bruijn terms requires shifting of indices, which complicates the rewrite rules or requires to introduce index shifting as a term former akin to explicit substitution.
Then, traditional e-graphs had no notion of binding, or variables in general: all e-graphs technically represent ground terms, with variables encoded as constants at the leaves of the graph.
Explicit substitution is then the preferred way of performing substitution of e-classes.
The alternative is to define substitution as a meta-operation on e-graphs, which requires the traversal of the e-graph, iterating through e-classes to see if they contain the required bound variable.
This involves obligations to make sure that such operation does not break the e-graph invariants.

The recent work of Schneider et al.~\cite{slotted-egraphs} addressed these issues in their \emph{slotted e-graphs} framework that introduced variables as first-class elements.
In slotted e-graphs, each e-class is parameterised by a set of free variables it uses and each e-node may invoke different instances of such an e-class by supplying concrete variable names for slots.
Consider the slotted e-graph example below.
 \[
  \adjustbox{scale=0.6}{
			\tikzfig{./figures/introduction/slotted-e-graph-example}
  }~.
\]
E-classes are represented as boxes with headers indicating the free variables they use.
There is also a single e-class used to represent the source of all varables: other variables are instantiated by invoking this e-class with different variable names.
This is an intentional design decision to enable sharing of terms that have the same \emph{free} variables up-to renaming.
For example, in a term $f(x,y) + f(z,w)$, with $x,y,z,w$ all free variables, the occurrence of $f$ is shared and exposes two generic free variables; the e-node $+$ then invokes the e-class containing $f$ with variables $x,y$ and $z,w$.
E-nodes are matched to their corresponding e-classes by computing a normal form with respect to the free variables, by canonically renaming them.
Note that e-class containing $2$ does not expose any variables and the e-class containing $\lambda x$ exposes one less variable than its child e-class.
Since variables are first-class citizens, slotted e-graphs have a substitution operation built-in that traverses the e-graph to perform the required substitution.
Yet, the practically preferred way of performing substitution is still to use explicit substitution.

We argue that string diagrams for $\catname{SLat}$-enriched closed symmetric monoidal categories, provide a more principled support for variable binding and substitution.
The syntax of such string diagrams is given in~\cref{fig:egraph-strings-closed}.
\begin{figure}
	\centering
	\adjustbox{scale=0.65}{
		\tikzfig{./figures/lambda/egraph-strings-2}
	}
	\caption{String diagrams for closed symmetric monoidal $\catname{SLat}$-categories. }
	\label{fig:egraph-strings-closed}
\end{figure}
String diagrams keep variables and e-classes (dashed boxes) as separate entities, variables are just wires connecting the diagram.
Substitution then does not require any special treatment, it is still just a rewiring of the diagram, regardless of whether we are connecting dashed boxes or not.
Finally, as we noted above, string diagrams have some notational advantages.

Consider the example of non-destructive $\beta$-reduction as performed in an e-graph and in a string diagram, in~\cref{fig:e-graph-substitution-string-combined}.
As noted, in the e-graph, explicit substitution e-nodes are added and then recursively resolved using additional rewrite rules.
This requires several steps and leaves explicit substitution nodes which are redundant when the substitution is done.
In contrast, in the string diagram, $\beta$-reduction is performed by a single rewrite step; the other steps perform structural $\catname{SLat}$-transformations to make the correspondence between the two structures more evident.
The corresponding slotted e-graph diagram may look as the one of traditional e-graphs if explicit substitution is used, or as string diagrammatic version, if the internal traversal is used to perform the substitution.
\begin{figure}
  \begin{subfigure}{\textwidth}
    \[
    \adjustbox{width=\linewidth}{
    \tikzfig{./figures/introduction/e-graph-explicit-substitution}
    }
    \]
    \subcaption{Explicit substitution in an e-graph.}
    \label{fig:explicit-substitution-egg-lambda-chapter}
\end{subfigure}
\begin{subfigure}{\textwidth}
    \[
    \adjustbox{width=\linewidth}{
    \tikzfig{./figures/introduction/e-string-substitution-2}
    }
    \]
    \subcaption{$\beta$-reduction in string diagrams.}
    \label{fig:beta-reduction-string-diagrams-lambda-chapter}    
\end{subfigure}
\caption{Performing $\beta$-reduction in an e-graph and in a string diagram.}
\label{fig:e-graph-substitution-string-combined}
\end{figure}

Yet, explicit substitution might be unavoidable when one wants to represent, for example, the equational theory of $\lambda_{\text{lsub}}^{\sim}$.
Moreover, for $\lambda_{\text{lsub}}^{\sim}$, the explicit substitution must coexist with the meta-substitution: it is the latter that performs the action under $ls$ rule.
We perform a little experiment to see the effect of graphical equivalence $\sim$ when performing equality saturation using slotted e-graphs.

\subsection{Modelling $\lambda_{\text{lsub}}^{\sim}$-calculus in slotted e-graphs}%
\label{sec:e-graph-explicit-subst}

% \begin{table*}
% 	\begin{tabular}{lcc}
% 		\texttt{beta}                  & $(\lambda x . t) u = t[x / u]$                       &                                                                    \\
% 		\texttt{subst\_same\_var}      & $x[x / u] = u$                                       &                                                                    \\
% 		\texttt{subst\_diff\_var}      & $y[x / u] = y$                                       & if $x \not \in \mathcal{F}(u)$                                     \\
% 		\texttt{subst\_application}    & $(uu')[x / v] = u[x / v]u'[x / v]$                   &                                                                    \\
% 		\texttt{subst\_lambda}         & $(\lambda x . u)[y / v] = \lambda x . (u[y / v])$    & if $x \not \in \mathcal{F}(v)$                                     \\
% 		\texttt{subst\_interchange\_1} & $(u[x / v])[x' / v'] = (u[x' / v'])[x / v[x' / v']]$ & if $x' \in \mathcal{F}(v)$ and $x \in \mathcal{F}(v')$             \\
% 		\texttt{subst\_interchange\_2} & $(u[x / v])[x' / v'] = (u[x / v[x' / v']])$          & if $x' \not \in \mathcal{F}(u)$                                    \\
% 		\texttt{subst\_comm}           & $(u[x / v])[x' / v'] = (u[x' / v'])[x / v]$          & if $x \not \in \mathcal{F}(v')$ and $x' \not \in \mathcal{F}(v)~.$ \\
% 	\end{tabular}
% 	\caption{Encoding of string diagram equivalence and explicit substitution for $\lambda_{\text{lsub}}$}
% 	\label{tbl:e-graph-subst-rules}
% \end{table*}

\begin{listing}
	\begin{minted}{rust}
pub enum Lambda {
   Lam(Bind<AppliedId>) = "lam",
   App(AppliedId, AppliedId) = "app",
   // subst \$x t u := t[x / u]
   Subst(Bind<AppliedId>, AppliedId) = "subst",
   Var(Slot) = "var",
}
	\end{minted}
	\caption{$\lambda_{\text{lsub}}$ syntax definition in slotted e-graphs}
	\label{listing:define_language}
\end{listing}


\begin{listing}
	\begin{minted}{rust}
fn beta() -> Rewrite<Lambda> {
    let pat = "(app (lam $1 ?b) ?t)";
    let outpat = "(subst $1 ?b ?t)";
    Rewrite::new("beta", pat, outpat)
}
	\end{minted}
	\caption{$\lambda_{\text{lsub}}$ $d\beta$-rule}
	\label{listing:lsub_beta_rule}
\end{listing}



Recall the grammar~\eqref{eq:lambda_lsub_grammar}, the reduction rules~\eqref{eq:lambda_lsub_reduction} and the equational theory given by equivalences~\eqref{eq:non_interacting_subst}--\eqref{eq:subst-application} of $\lambda_{\text{lsub}}^{\sim}$.
These can be directly encoded in the framework of slotted e-graphs of Schneider et al.~\cite{slotted-egraphs} as can be seen in \cref{listing:define_language}, \cref{listing:lsub_beta_rule}, \cref{listing:lsub_ls_rules-1}, \cref{listing:lsub_ls_rules-2}, \cref{listing:gc_rule}, and \cref{listing:graphical_equivalence}.
We use \mintinline{rust}{"subst"} to denote explicit substitution and use the built-in substitution mechanism provided by slotted e-graphs to encode $C\llbracket x \rrbracket[x / u] \to C\llbracket u \rrbracket[x/u]$ rules.
Notably, since equivalence is symmetric, the corresponding rules have two versions---\textit{left} and \textit{right}.
We check for the side conditions using conditional rewrites provided by the slotted e-graphs framework.

To show the effect of those particular sets of rewrites when performing equality saturation for $\lambda_{\text{lsub}}^{\sim}$ we apply certain subsets of these rewrites to a simple term that encodes $2 + 2$ using Church encoding.
When doing so, we keep track of the number of e-nodes in the saturated e-graph as well as the number of times each rewrite was applied during saturation (we count only the applications that change the e-graph).
We also note whether a particular rewrite is absorbed by using string-diagrammatic (as a concrete combinatorial object) representation of $\lambda_{\text{lsub}}^{\sim}$-terms.
The idea is to show how much workload can be absorbed by using the above-mentioned representation that represents certain equivalences on the nose.
\Cref{tbl:slotted} shows the number of e-nodes after saturation with different sets of rewrite rules associated with $\lambda_{\text{lsub}}^{\sim}$ enabled.
\texttt{subst\_lambda} is the distributivity of substitution over lambda abstraction (encoded using two rewrite rules) corresponding to equivalence~\eqref{eq:subst-lambda}; \texttt{subst\_comm} is the commutativity of substitutions which corresponds to the equivalence~\eqref{eq:non_interacting_subst}, whereas \texttt{app\_subst} corresponds to the equivalence~\eqref{eq:subst-application}.
We can see that enabling all those rules---that is truly performing rewriting modulo graphical equivalence---substantially increases the number of e-nodes in the e-graph.
Since \texttt{subst\_comm} and \texttt{app\_subst} are absorbed by the string-diagrammatic representation we theoretically can get away with fewer nodes if we used this representation.

Another metric that we use is the number of times each rewrite was applied during saturation which is shown in \cref{tbl:slotted-rewrite-applications}.
$gc$ rule is encoded using \texttt{subst\_no\_var} rewrite rule.
We can see that the rules that are absorbed by the string-diagrammatic representation (highlighted in grey) constitute about $40\%$ of all the rewrite applications during saturation when all the rules are enabled.
So again, using a more suitable representation can reduce the workload of the equality saturation in this case.

This particular calculus of choice might seem rather exotic.
Indeed, $\lambda_{\text{lsub}}^{\sim}$ is designed so that reduction mirrors cut elimination in proof nets and terms are quotiented, so that different terms under $\sim$ correspond to different proof nets.
However, the phenomenon illustrated here is more general: non-interacting \texttt{let}-bindings are not exotic and represent a reasonable notion of equivalence on programs: the order of \texttt{let}-bindings must be immaterial as long as they do not interfere with each other.

% \begin{remark}
% Even though $gc$ rule is not absorbed by the e-hypergraph representation (despite being a part of string diagram equivalence), we can argue that this rule can be disregarded when using e-hypergraphs.
% Consider a substitution $y[x / v]$ which corresponds to the following composition in the internal language
% \[
% \inferrule*{\Gamma_1,y : T, \Gamma_2 \vdash v : T' \qquad \Gamma_1, x : T', y : T, \Gamma_2 \vdash y : T}{
% 	\Gamma_1, y : T, \Gamma_2 \vdash y[x / v] : T \equiv y
% }
% \]
% This composition is represented using the string diagram in~\autoref{fig:subst_no_var_string}.
% While string diagrammatically this composition equals the string diagram for just $y$ since we are working modulo the equations for copy and delete, they are not isomorphic as e-hypergraphs.
% However, since the output of $v$ gets deleted anyway, we can ignore any rewrites that involve $v$ and its inputs.
% The deletion morphisms for $\Gamma_1, \Gamma_2$ that follow the copy can be absorbed by the e-hypergraph representation by following the approach of Milosavljevi\'{c} et al.~\cite{zanassi_comonoid}.
% \end{remark}

\begin{listing}
	\begin{minted}{rust}
fn subst_no_var() -> Rewrite<Lambda> {
   let pat = "(subst $x ?t1 ?t2)";
   let outpat = "?t1";
   Rewrite::new_if("subst_no_var", pat, outpat, |subst, _ | {
      !subst["t1"].slots().contains(&Slot::named("x"))
   })
}
	\end{minted}
	\caption{$\lambda_{\text{lsub}}$ $gc$-rule}
	\label{listing:gc_rule}
\end{listing}


% \begin{figure}
% 	\centering
% 	\tikzfig{./figures/subst_no_var_string}
% 	\caption{String diagram for $y[x/v]$}
% 	\label{fig:subst_no_var_string}
% \end{figure}

\begin{table*}
	\begin{tabular}{lcccc}
		Set of equations & $d\beta$ + $ls$ + $gc$ & $d\beta$ + $ls$ + $gc$    & $d\beta$ + $ls$ + $gc$     & $d\beta$ + $ls$ + $gc$\\
		                 &                        & + \texttt{subst\_lambda}  & + \texttt{subst\_lambda}   & + \texttt{subst\_lambda}\\
		                 &                        &                           & + \texttt{subst\_comm}     & + \texttt{subst\_comm}\\
						 &                        &                           &                            & + \texttt{subst\_app}\\						 
		\# of e-nodes    & 88                     & 185                       & 244                        & 846
	\end{tabular}
	\caption{Number of e-nodes after saturating the e-graph for $2 + 2$ with different sets of rewrite rules for $\lambda_{\text{lsub}}$}
	\label{tbl:slotted}
\end{table*}

\begin{table*}
  \adjustbox{max width=\textwidth}{
	\begin{tabular}{l|ccccc}
		Rewrite rule & \# of applications     & \# of applications         & \# of applications      & \# of applications       & Absorbed?\\
		             &($d\beta$ + $ls$ + $gc$)&($d\beta$ + $ls$ + $gc$     &($d\beta$ + $ls$ + $gc$  &($d\beta$ + $ls$ + $gc$   &\\
		             &                        & + \texttt{subst\_lambda}  & + \texttt{subst\_lambda}& + \texttt{subst\_lambda} &\\
		             &                        &                            & + \texttt{subst\_comm} & + \texttt{subst\_comm}   &\\
		             &                        &                            &                         & + \texttt{subst\_app})   &\\
		\hline
		\texttt{beta} & 14 & 57 & 49 & 80 & No\\
		\texttt{subst\_no\_var} & 35 & 113 & 118 & 197 & No\\
		% \rowcolor{gray!30}
		\texttt{ls\_rule\_*} & 48 & 141 & 140 & 302 & No\\
		\texttt{subst\_lambda} & --- & 76 & 109 & 336 & No\\
		\rowcolor{gray!30}
		\texttt{subst\_comm} & --- & --- & 104 & 321 & Yes\\
		\rowcolor{gray!30}
		\texttt{subst\_app} & --- & --- & --- & 263 & Yes\\
	\end{tabular}}
	\caption{Number of times each rewrite rule was applied during saturation of the e-graph for $2 + 2$}
	\label{tbl:slotted-rewrite-applications}
\end{table*}

% The \enquote{explicit subst} set of rewrite rules consist of the first five rules of~\autoref{tbl:e-graph-subst-rules}.
% We can see that additionally saturating the e-graph with respect to interchange and commutativity rules increases the number of e-nodes by a factor of 2.
% These are the exact equations that are absorbed by using string-diagrammatic representation of $\lambda_{\text{lsub}}$-terms.
% Notably, string diagrams (and (closed e-)hypergraphs) do not absorb the equation \texttt{subst\_lambda} but we believe it is practically feasible to do rewriting modulo this rule by traversing through the boxes when searching for a redex.
% Furthermore, (cospans of) closed e-hypergraphs, as presented in~\autoref{sec:closed-e-hyp}, do not absorb the equations for associativity, commutativity and unitality for $\comonoid$ and $\counit$, but the constraints on the cospans and the boundary complement may be relaxed to absorb those as well by following the work of~\cite{zanassi_comonoid} with which our approach is completely compatible.


\clearpage

\begin{listing*}
	\centering
\begin{minipage}{0.8\textwidth}
	\begin{minted}{rust}
fn ls_rule_var() -> Rewrite<Lambda> {
   let pat = "(subst $x (var $x) ?t)";
   let outpat = "(subst $x (var $x)[(var $x) := ?t] ?t)";
   Rewrite::new("ls_var", pat, outpat)
}

fn ls_rule_app_l() -> Rewrite<Lambda> {
   let pat = "(subst $x (app ?u ?v) ?t)";
   let outpat = "(subst $x (app ?u[(var $x) := ?t] ?v) ?t)";
   Rewrite::new_if("ls_app_l", pat, outpat, |subst, _| {
      subst["u"].slots().contains(&Slot::named("x")) 
	  && !subst["v"].slots().contains(&Slot::named("x"))
    })
}

fn ls_rule_app_r() -> Rewrite<Lambda> {
   let pat = "(subst $x (app ?u ?v) ?t)";
   let outpat = "(subst $x (app ?u ?v[(var $x) := ?t]) ?t)";
   Rewrite::new_if("ls_app_r", pat, outpat, |subst, _| {
      !subst["u"].slots().contains(&Slot::named("x")) 
	  && subst["v"].slots().contains(&Slot::named("x"))
   })
}
	\end{minted}
\end{minipage}
	\caption{$\lambda_{\text{lsub}}$ $ls$-rules (part 1).}
	\label{listing:lsub_ls_rules-1}
\end{listing*}

\begin{listing*}
	\centering
\begin{minipage}{0.8\textwidth}
	\begin{minted}{rust}
fn ls_rule_lam() -> Rewrite<Lambda> {
   let pat = "(subst $x (lam $y ?b) ?t)";
   let outpat = "(subst $x (lam $y ?b[(var $x) := ?t]) ?t)";
   Rewrite::new_if("ls_lam", pat, outpat, |subst, _| {
      subst["b"].slots().contains(&Slot::named("x"))
   })
}

fn ls_rule_subst_l() -> Rewrite<Lambda> {
   let pat = "(subst $x (subst $y ?u ?t1) ?t2)";
   let outpat = "(subst $x (subst $y ?u ?t1[(var $x) := ?t2]) ?t2)";
   Rewrite::new_if("ls_subst_l", pat, outpat, |subst, _| {
      subst["t1"].slots().contains(&Slot::named("x"))
   })
}

fn ls_rule_subst_r() -> Rewrite<Lambda> {
    let pat = "(subst $x (subst $y ?u ?t1) ?t2)";
    let outpat = "(subst $x (subst $y ?u[(var $x) := ?t2] ?t1) ?t2)";
    Rewrite::new_if("ls_subst_r", pat, outpat, |subst, _| {
          subst["u"].slots().contains(&Slot::named("x"))
    })
}
	\end{minted}
\end{minipage}
	\caption{$\lambda_{\text{lsub}}$ $ls$-rules (part 2).}
	\label{listing:lsub_ls_rules-2}
\end{listing*}

\clearpage

\begin{listing*}
	\centering
	\begin{minipage}{0.8\textwidth}
	\begin{minted}{rust}
fn subst_comm() -> Rewrite<Lambda> {
   let pat = "(subst $y (subst $x ?u ?t1) ?t2)";
   let outpat = "(subst $x (subst $y ?u ?t2) ?t1)";
   Rewrite::new_if("subst_comm", pat, outpat, |subst,_| {
      !subst["t2"].slots().contains(&Slot::named("x")) 
      && !subst["t1"].slots().contains(&Slot::named("y"))
   })
}

fn subst_lambda_l() -> Rewrite<Lambda> {
    let pat = "(subst $x (lam $y ?b) ?t)";
    let outpat = "(lam $y (subst $x ?b ?t))";
    Rewrite::new_if("subst_lambda_l", pat, outpat, |subst, _| {
        !subst["t"].slots().contains(&Slot::named("y"))
    })
}

fn subst_lambda_r() -> Rewrite<Lambda> {
   let pat = "(lam $y (subst $x ?b ?t))";
   let outpat = "(subst $x (lam $y ?b) ?t)";
   Rewrite::new_if("subst_lambda_r", pat, outpat, |subst, _| {
      !subst["t"].slots().contains(&Slot::named("y"))
   })
}

fn subst_app_l() -> Rewrite<Lambda> {
   let pat = "(subst $x (app ?t1 ?t2) ?t3)";
   let outpat = "(app (subst $x ?t1 ?t3) ?t2)";
   Rewrite::new_if("subst_app_2", pat, outpat, |subst,_| {
      !subst["t2"].slots().contains(&Slot::named("x"))
   })
}
fn subst_app_r() -> Rewrite<Lambda> {
   let pat = "(app (subst $x ?t1 ?t3) ?t2)";
   let outpat = "(subst $x (app ?t1 ?t2) ?t3)";
   Rewrite::new_if("subst_app_2_2", pat, outpat, |subst,_| {
      !subst["t2"].slots().contains(&Slot::named("x"))
   })
}
	\end{minted}
\end{minipage}
	\caption{Graphical equivalence rules for $\lambda_{\text{lsub}}$.}
	\label{listing:graphical_equivalence}
\end{listing*}

The preceding example motivates moving part of the equational theory
from the rewrite system into the representation itself. 
We now make that representation precise. 
We first define a typed raw syntax with
abstraction and finite joins, then introduce the corresponding closed
e-hypergraphs and interpret the syntax into them.
We follow the same approach as in~\cref{chapter:monoidal}.
In particular, the objects we construct absorb the symmetric monoidal structure on terms, but they can further be looked at though the prism of $\lambda$-calculus: $\beta$-reduction becomes a simple rewrite rule, vertices naturally represent substitution and binding and the hierarchical structure supports equivalent terms.

\section{Raw syntax}
We start with the notion of closed-monoidal terms.

\begin{definition}[Objects]
Let $\Sigma=(\Sigma_0,\Sigma_1)$ be a monoidal signature.
The set $\operatorname{Ob}^{\Lambda}_{\Sigma}$ consists of
object expressions generated by
\[
  A,B ::= I \mid a \in \Sigma_0
          \mid A \otimes B
          \mid A \multimap B,
\]
modulo the least congruence generated by
\[
  (A \otimes B) \otimes C = A \otimes (B \otimes C),
  \qquad
  I \otimes A = A = A \otimes I.
\]
No other equations are imposed on objects. Thus $\otimes$ is strictly
associative and unital, while $\multimap$ is a free type constructor.

A word
\[
  \langle a_1,\ldots,a_n\rangle\in\Sigma_0^*
\]
is regarded as the object
\[
  a_1\otimes\cdots\otimes a_n,
\]
with the empty word \(\varepsilon=\langle\,\rangle\) interpreted as
\(I\). In particular, the domains and codomains of generators in
\(\Sigma_1\) determine objects of
\(\operatorname{Ob}^{\Lambda}_{\Sigma}\).

An object is \emph{atomic} if it is a base type \(a\in\Sigma_0\) or an
internal-hom type \(A\multimap B\). Write
\(\operatorname{At}^{\Lambda}_{\Sigma}\) for the set of atomic types.
Every object has a unique expression as a finite tensor word of atomic
types.

We display words of atomic types as lists
\[
  \langle T_1,\ldots,T_n\rangle
  \in
  \bigl(\operatorname{At}^{\Lambda}_{\Sigma}\bigr)^*,
\]
and use juxtaposition for concatenation of words. Define
\[
  \operatorname{wd}:
  \operatorname{Ob}^{\Lambda}_{\Sigma}
  \longrightarrow
  \bigl(\operatorname{At}^{\Lambda}_{\Sigma}\bigr)^*
\]
by
\[
  \operatorname{wd}(I)=\varepsilon,
  \qquad
  \operatorname{wd}(A\otimes B)
  =
  \operatorname{wd}(A)\operatorname{wd}(B),
  \qquad
  \operatorname{wd}(A\multimap B)
  =
  \langle A\multimap B\rangle.
\]

Conversely, define the object represented by a word by
\[
  \left\lVert\langle T_1,\ldots,T_n\rangle\right\rVert_{\otimes}
  \defeq
  T_1\otimes\cdots\otimes T_n,
  \qquad
  \lVert\varepsilon\rVert_{\otimes}\defeq I.
\]
The functions \(\operatorname{wd}\) and
\(\lVert-\rVert_{\otimes}\) are mutually inverse.
\end{definition}

\begin{example}[Words and objects]
Let \(a,b,c,d\in\Sigma_0\). The expression
\(a\multimap(b\otimes c)\) is one atomic type, and
\[
  \operatorname{wd}
  \bigl((a\multimap(b\otimes c))\otimes a\otimes d\bigr)
  =
  \langle a\multimap(b\otimes c),a,d\rangle.
\]
This three-letter word represents the object
\[
  (a\multimap(b\otimes c))\otimes a\otimes d.
\]
By contrast, \((a\multimap b)\otimes c\) is represented by the
two-letter word
\[
  \langle a\multimap b,c\rangle.
\]
\end{example}

\begin{definition}[Raw closed-monoidal terms]
For $A,B \in \operatorname{Ob}^{\Lambda}_{\Sigma}$, the
families $\operatorname{Tm}^{\Lambda}_{\Sigma}(A,B)$ of
\emph{raw closed-monoidal terms} are the least families
generated by the following constructors:
\begin{enumerate}
  \item Each generator $f \in \Sigma_1$ is a term
  \[
    f : \mathrm{dom}(f) \to \mathrm{cod}(f),
  \]
  where its domain and codomain words are regarded as objects.

  \item There are primitive identity terms
  \[
    \id_I : I \to I,
    \qquad
    \id_T : T \to T
    \quad
    (T \in \operatorname{At}^{\Lambda}_{\Sigma}).
  \]

  \item For all $T,U \in \operatorname{At}^{\Lambda}_{\Sigma}$,
  there is a primitive symmetry term
  \[
    \sym_{T,U} : T \otimes U \to U \otimes T.
  \]

  \item If $f : A \to B$ and $g : B \to C$ are terms,
  then
  \[
    f \fatsemi g : A \to C
  \]
  is a term.

  \item If $f : A \to B$ and $g : C \to D$ are terms,
  then
  \[
    f \otimes g : A \otimes C \to B \otimes D
  \]
  is a term.

  \item For all $A \in \operatorname{At}^{\Lambda}_{\Sigma}$ and $B \in \operatorname{Ob}^{\Lambda}_{\Sigma}$,
  there is an evaluation term
  \[
    \ev_{A,B} : (A \multimap B) \otimes A \to B.
  \]

  \item For all $X \in \operatorname{At}^{\Lambda}_{\Sigma}$ and $A,B \in \operatorname{Ob}^{\Lambda}_{\Sigma}$, if $f : A \otimes X \to B$ is a term, then
  \[
    \Lambda_{A,X,B}(f) : A \to X \multimap B
  \]
  is a term.
\end{enumerate}
\end{definition}

\begin{remark}
  Above we assume that evaluation and abstraction happen with respect to atomic types.
  This is to simplify the necessary bookkeeping in the combinatorial definitions.
\end{remark}

We write $\id_A$ and $\sym_{A,B}$ for the standard terms on arbitrary
objects derived from the primitive identities and symmetries using
$\fatsemi$ and $\otimes$.

\begin{definition}[Raw closed-monoidal terms with joins]
The families
$\operatorname{Tm}^{\Lambda,\join}_{\Sigma}(A,B)$
are generated by the same constructors as
$\operatorname{Tm}^{\Lambda}_{\Sigma}(A,B)$, together with
the following constructor:
\[
  \frac{f_1,\ldots,f_n : A \to B \qquad (n \geq 2)}
       {f_1 \join \cdots \join f_n : A \to B}.
\]
\end{definition}
\begin{remark}
In this chapter we are concerned with raw $\Sigma^{\join}$-terms of $\operatorname{Tm}^{\Lambda, \join}_{\Sigma}$, rather than with morphisms as was the case with $\mathbf{S}^{\join}_{\Sigma}$ and $\mathbf{H}^{\join}_{\Sigma}$.
Because string diagrams have $n$-ary boxes, we also need $n$-ary join at the level of raw syntax, to make the eventual correspondence between terms and combinatorial structures precise.
\end{remark}

We shall show that rewriting terms in
$\operatorname{Tm}_{\Sigma}^{\Lambda,\join}$ modulo the symmetric monoidal
and $\catname{SLat}$ equations can equivalently be performed on their
combinatorial representation in terms of \emph{closed e-hypergraphs}.
Recall that string diagrams as in~\cref{sec:string-diagrams-Bonchi} are defined as terms quotiented by the axioms of a symmetric monoidal category.
Thus, the set of terms $\operatorname{Tm}^{\Lambda, \join}_{\Sigma}$ defines exactly string diagrams for $\catname{SLat}$-enriched symmetric monoidal category. 

\begin{definition}[Structural congruences]
\label{def:lambda-structural-congruences}
Let $\equiv_0$ be the congruence on
$\operatorname{Tm}^{\Lambda}_{\Sigma}$ generated by all well-typed
instances of the strict symmetric monoidal equations in
\cref{fig:string-diagram-axioms}.

Let $\equiv_{\mathsf{str}}$ be the congruence on
$\operatorname{Tm}^{\Lambda,\join}_{\Sigma}$ generated by the following
well-typed equations:
\begin{enumerate}
  \item The strict symmetric monoidal equations, allowing terms that
  contain joins.

  \item The semilattice equations. For $n\geq2$ and parallel terms
  $f_1,\ldots,f_n:A\to B$, every permutation $\theta$ of $[n]$ gives
  \[
    f_1 \join \cdots \join f_n
    \equiv_{\mathsf{str}}
    f_{\theta(1)} \join \cdots \join f_{\theta(n)},
  \]
  a join nested in one argument may be flattened,
  \begin{align*}
    f_1 \join \cdots \join &f_{i-1}
    \join (g_1 \join \cdots \join g_m)
    \join f_{i+1} \join \cdots \join f_n
    \equiv_{\mathsf{str}}\\
    &f_1 \join \cdots \join f_{i-1}
    \join g_1 \join \cdots \join g_m
    \join f_{i+1} \join \cdots \join f_n,
  \end{align*}
  and a repeated argument may be contracted: if $f_i=f_j$ for some
  $i\neq j$, then
  \[
    f_1 \join \cdots \join f_n
    \equiv_{\mathsf{str}}
    f_1 \join \cdots \join f_{j-1}
    \join f_{j+1} \join \cdots \join f_n .
  \]
  We identify a one-element join with the term itself, so that contracting a binary join gives
  $f\join f\equiv_{\mathsf{str}}f$. Together, these are the $n$-ary
  forms of commutativity, associativity and idempotence. 

  \item Preservation of joins in each argument of
  composition and tensor: for $n\geq2$,
  \begin{align*}
    (f_1 \join \cdots \join f_n) \fatsemi h
      &\equiv_{\mathsf{str}}
      (f_1 \fatsemi h) \join \cdots \join (f_n \fatsemi h),\\
    h \fatsemi (f_1 \join \cdots \join f_n)
      &\equiv_{\mathsf{str}}
      (h \fatsemi f_1) \join \cdots \join (h \fatsemi f_n),\\
    (f_1 \join \cdots \join f_n) \otimes h
      &\equiv_{\mathsf{str}}
      (f_1 \otimes h) \join \cdots \join (f_n \otimes h),\\
    h \otimes (f_1 \join \cdots \join f_n)
      &\equiv_{\mathsf{str}}
      (h \otimes f_1) \join \cdots \join (h \otimes f_n),
  \end{align*}
  \item Preservation of joins by abstraction: for $n\geq2$,
  \begin{equation}
    \Lambda_{X,A,B}(f_1 \join \cdots \join f_n)
    \equiv_{\mathsf{str}}
    \Lambda_{X,A,B}(f_1) \join \cdots \join \Lambda_{X,A,B}(f_n),
    \label{eq:lambda-join-structural}
  \end{equation}
  for all $f_1,\ldots,f_n : X \otimes A \to B$.
\end{enumerate}
Items~2 and~3 are precisely the equations of \cref{fig:string-equations},
read as equations between raw closed-monoidal terms; item~4 is the additional equation
for abstraction.
\end{definition}

Note that the structural congruences of \cref{def:lambda-structural-congruences} mention neither of equations~\eqref{eq:lambda_1}--\eqref{eq:lambda_3}.
Quotienting with respect to these equations would result in an $\catname{SLat}$-enriched closed symmetric monoidal category, so that rewriting would have to be performed moduloe closed structure as well.
As these equations have operational significance, we do not apply that additional quotient, and strictly speaking, consider just $\catname{SLat}$-enriched symmetric monoidal categories, altough throughout this chapter we just consider raw closed-monoidal terms.

% \subsection*{Categorical foundations of slotted e-graphs}
% We take a brief pause to discuss the equation~\eqref{eq:lambda-join-structural} with respect to slotted e-graphs of Schneider et al.~\cite{slotted-egraphs}.
% Slotted e-graphs extend the traditional signature employed by e-graphs with an explicit construct for binding ($\text{bind}\; \$x\; tc$) and variables (\emph{slots}).
% 	Their syntax is generated by the following grammar:
% 	\begin{alignat*}{2}
% 		\text{(functional symbols)} &        &                            & \; f,\;g, \ldots ;                                                                       \\
% 		\text{(slots)}              &        &                            & \; \$x,\; \$y, \ldots ;                                                                   \\
% 		\text{(term children)}      &        & tc               \Coloneqq & \; t \mid tc_{1},\; \ldots,\; tc_{n} \mid \text{bind}\; \$x\; tc \mid \$x ; \qquad n \geq 1 \\
% 		\text{(terms)}              & \qquad & t                \Coloneqq & \; f \mid f(tc) .
% 	\end{alignat*}
% 	The traditional congruence relation maintained by e-graphs for a set of equations $E$ is given below (plus the obvious rules for transitivity, symmetry and reflexivity of $\cong$; in this subsection only, $\cong$ denotes the e-graph congruence rather than isomorphism):
% 	\begin{mathpar}
% 		\inferrule*[right=start]
% 		{a = b \in E}
% 		{E \vdash a \cong b}
% 		\quad
% 		\inferrule*[right=cong. (variadic)]
% 		{E \vdash tc_i \cong tc'_i,\; i = 1, \ldots, k}
% 		{E \vdash tc_1, \ldots, tc_k \cong tc'_1, \ldots, tc'_k}
% 		\quad 
% 		\inferrule*[right=cong. (f)]
% 		{E \vdash tc \cong tc'}
% 		{E \vdash f(tc) \cong f(tc')}
% 	\end{mathpar}
%   Recall that we interpret the e-graph congruence by putting the equivalent terms together with a $\join$.
%   The rule $tc \cong tc' \implies f(tc) \cong f(tc')$ is then exactly the property of composition in an $\catname{SLat}$-category, namely $f(tc \join tc') = f(tc) \join f(tc')$.
% 	Slotted e-graphs extend the above rules with the following:
% 	\begin{equation}%
% 		\label{eq:cong-bind}
% 		\inferrule*[right=cong. (bind)]
% 		{E \vdash tc \cong tc'}
% 		{E \vdash \text{bind}\; \$x\; tc \cong \text{bind}\; \$x\; tc'},
% 	\end{equation}
% 	plus the rules for equivalence of terms up to renaming of slots.
%   This rule is exactly the property of preservation of joins by abstraction in~\eqref{eq:lambda-join-structural}.
%   We did not have to introduce this rule in an ad hoc manner, nor did we have to justify it---it is provided by the axioms of the category we work in and thus the rule is sound with respect to semantics.

\section{Closed e-hypergraphs}
Let $\Gamma_\Sigma$ be the typed monoidal signature whose object
types are $\operatorname{At}^{\Lambda}_{\Sigma}$ and whose generators
are
\[
  (\Gamma_\Sigma)_1
  =
  \Sigma_1
  \sqcup
  \{
    \ev_{A,B}
    \mid
    A \in \operatorname{At}^{\Lambda}_{\Sigma},B\in\operatorname{Ob}^{\Lambda}_{\Sigma}
  \}.
\]
A generator $f\in\Sigma_1$ retains its original domain and codomain,
regarded as words of atomic types, while
\[
  \ev_{A,B}:
  \operatorname{wd}(A\multimap B)\operatorname{wd}(A)
  \longrightarrow
  \operatorname{wd}(B).
\]

\begin{definition}[Closed rooted e-hypergraph]
\label{def:lambda-rooted-graph}
A \emph{closed rooted e-hypergraph} over \(\Sigma\) is a tuple
\[
  \mathcal G
  =
  (V,E,s,t,l_V,l_E,\top,<^\mu)
\]
with the same underlying hypergraph and hierarchical data as a rooted
e-hypergraph of \cref{def:rooted_ehypergraph}, except for the following
changes.

Vertices and edges are labelled by
\[
  l_V:
  V\longrightarrow\operatorname{At}^{\Lambda}_{\Sigma},
  \qquad
  l_E:
  E\longrightarrow
  (\Gamma_\Sigma)_1
  \sqcup
  \{\boxsort,\cmpsort,\lambdasort\}.
\]
The vertex labelling extends pointwise to words:
\[
  l_V^*:
  V^*
  \longrightarrow
  \bigl(\operatorname{At}^{\Lambda}_{\Sigma}\bigr)^*,
  \qquad
  l_V^*\bigl(\langle v_1,\ldots,v_n\rangle\bigr)
  =
  \langle l_V(v_1),\ldots,l_V(v_n)\rangle.
\]

Write
\[
  E_{\Gamma_\Sigma},
  \qquad
  E_{\boxsortsubscript},
  \qquad
  E_{\cmpsortsubscript},
  \qquad
  E_{\lambdasortsubscript}
\]
for the respective inverse images under \(l_E\).
Edges in \(E_{\Gamma_\Sigma}\) are called \emph{ordinary edges}, and
edges in \(E_{\lambdasortsubscript}\) are called
\emph{\(\lambda\)-edges}.

The tree-structure condition of
\cref{def:rooted_ehypergraph} is unchanged. 
As before, for \(x\neq\top\) we write
\(<^\mu(x)\), or \(<^\mu_{\mathcal G}(x)\), for the immediate parent of
\(x\). The remaining conditions are replaced by the following.

\begin{enumerate}
  \item \emph{Typing.}
  Every ordinary edge \(e\in E_{\Gamma_\Sigma}\) satisfies
  \[
    l_V^*(s(e))
    =
    \operatorname{dom}_{\Gamma_\Sigma}(l_E(e)),
    \qquad
    l_V^*(t(e))
    =
    \operatorname{cod}_{\Gamma_\Sigma}(l_E(e)).
  \]
  Every component edge \(c\in E_{\cmpsortsubscript}\) satisfies
  \[
    s(c)=t(c)=\varepsilon.
  \]

  \item \emph{Alternation.}
  The immediate parent of every vertex, ordinary edge,
  \(\boxsort\)-edge and \(\lambdasort\)-edge is a component edge.
  The immediate parent of every component edge other than \(\top\)
  is either a \(\boxsort\)-edge or a \(\lambdasort\)-edge.

  \item \emph{Incidence.}
  If
  \[
    v\in s(e)
    \quad\text{or}\quad
    v\in t(e)
  \]
  for
  \[
    e\in
    E_{\Gamma_\Sigma}
    \cup E_{\boxsortsubscript}
    \cup E_{\lambdasortsubscript},
  \]
  then \(v\) and \(e\) have the same immediate parent.

  \item \emph{Hierarchical arity.}
  Every \(\boxsort\)-edge has at least two component children, and
  every \(\lambdasort\)-edge has exactly one component child.
\end{enumerate}

As before, \(\mathfrak T\) denotes the closed rooted e-hypergraph with
no vertices and with its component root as its only edge.
\end{definition}

The only new hierarchical sort is $\lambdasort$. A $\lambda$-edge has one
component, representing the body of an abstraction. Evaluation edges are
ordinary edges labelled by the generators
$\ev_{A,B}\in(\Gamma_\Sigma)_1$.
These $\lambdasort$-edges were introduced by Ghica et al.~\cite{fscd} and served as the inspiration for our $\boxsort$-edges in~\cref{chapter:monoidal}.

\begin{remark}[Atomic wire types]
Each vertex represents one wire and is therefore labelled by one atomic
type. An arbitrary object $A$ is represented by the ordered list of
vertices labelled by $\operatorname{wd}(A)$. In particular,
$A\multimap B$ labels one vertex, whereas a non-trivial tensor is
represented by several vertices.
\end{remark}

\begin{example}[Extending vertex labels to words]
Let \(a,b,c,d\in\Sigma_0\), and suppose that
\[
  l_V(v_1)=a\multimap(b\otimes c),
  \qquad
  l_V(v_2)=a,
  \qquad
  l_V(v_3)=d.
\]
The pointwise extension of \(l_V\) gives
\[
  l_V^*\bigl(\langle v_1,v_2,v_3\rangle\bigr)
  =
  \langle a\multimap(b\otimes c),a,d\rangle.
\]
This is a three-letter word of atomic types. The object it represents is
\[
  \left\lVert
    \langle a\multimap(b\otimes c),a,d\rangle
  \right\rVert_{\otimes}
  =
  (a\multimap(b\otimes c))\otimes a\otimes d.
\]
\end{example}

\begin{example}
An example of a closed e-hypergraph is shown below.
\[
\adjustbox{scale=0.8}{
\tikzfig{./figures/lambda/closed_e_hypergraph_example}~.
}
\]
We depict $\lambdasort$-edges as rounded boxes. They are distinguished
from $\boxsort$-edges by having exactly one component child.
\end{example}

Closed rooted e-hypergraphs come with the familiar classes of homomorphisms.
\begin{definition}[Lax, anchored and strict homomorphisms]
\label{def:lambda-rooted-homomorphism}
Let $\mathcal F$ and $\mathcal G$ be closed rooted
e-hypergraphs over $\Sigma$.

A \emph{lax homomorphism}
\[
  \phi:\mathcal F\longrightarrow\mathcal G
\]
is a pair of functions
\[
  \phi_V:V_{\mathcal F}\longrightarrow V_{\mathcal G},
  \qquad
  \phi_E:E_{\mathcal F}\longrightarrow E_{\mathcal G}
\]
satisfying the incidence- and ancestry-preservation conditions
(1) and (3) of \cref{def:rooted_homomorphism}, and preserving both
labelling functions:
\[
  l_{V,\mathcal G}(\phi_V(v))=l_{V,\mathcal F}(v)
  \qquad
  (v\in V_{\mathcal F}),
\]
and
\[
  l_{E,\mathcal G}(\phi_E(e))=l_{E,\mathcal F}(e)
  \qquad
  (e\in E_{\mathcal F}).
\]

It is \emph{anchored} if it preserves immediate parents, as in
condition~(3$'$) of \cref{def:rooted_homomorphism}, and \emph{strict}
if it is anchored and
\[
  \phi_E(\top_{\mathcal F})=\top_{\mathcal G}.
\]
For an anchored homomorphism, the component edge
$\phi_E(\top_{\mathcal F})$ is called its \emph{anchor}.

Closed rooted e-hypergraphs with lax homomorphisms form a category
$\catname{CEHyp}_{\leq}(\Sigma)$. Anchored and strict homomorphisms
form wide subcategories
\[
  \catname{CEHyp}_{\mathrm{s}}(\Sigma)
  \subseteq
  \catname{CEHyp}_{\bullet}(\Sigma)
  \subseteq
  \catname{CEHyp}_{\leq}(\Sigma).
\]
\end{definition}

\subsection{Typed interfaces and extended cospans}

\begin{definition}[Typed finite interfaces]
\label{def:lambda-typed-interfaces}
The category \(\mathbb F^{\Lambda}_{\Sigma}\) has as objects pairs
\[
  U=([n],\tau),
  \qquad
  \tau:[n]\longrightarrow
  \operatorname{At}^{\Lambda}_{\Sigma},
\]
where \(\tau(i)\) is the atomic type assigned to the \(i\)-th interface
position.

A morphism
\[
  \varphi:([n],\tau)\longrightarrow([m],\upsilon)
\]
is a function \(\varphi:[n]\to[m]\) that preserves types:
\[
  \tau(i)=\upsilon(\varphi(i))
  \qquad
  (i\in[n]).
\]
Identities and composition are inherited from \(\catname{Set}\).

For an ordered subset
\[
  S=\{i_1<\cdots<i_k\}\subseteq[n],
\]
define its \emph{type word} by
\[
  \operatorname{ty}_{U}(S)
  \defeq
  \langle\tau(i_1),\ldots,\tau(i_k)\rangle,
  \qquad
  \operatorname{ty}_{U}(\varnothing)\defeq\varepsilon.
\]

The category \(\mathbb F^{\Lambda}_{\Sigma}\) is strict symmetric
monoidal. Its tensor product on objects is concatenation:
\[
  ([n],\tau) \coprod ([m],\upsilon)
  \defeq
  ([n+m],\tau + \upsilon),
\]
where
\[
  (\tau + \upsilon)(i)
  \defeq
  \begin{cases}
    \tau(i), & 1\leq i\leq n,\\
    \upsilon(i-n), & n<i\leq n+m.
  \end{cases}
\]

For morphisms
\[
  \varphi:([n],\tau)\to([n'],\tau'),
  \qquad
  \psi:([m],\upsilon)\to([m'],\upsilon'),
\]
their tensor product is the block sum
\[
  \varphi \coprod \psi:[n+m]\longrightarrow[n'+m']
\]
defined by
\[
  (\varphi \coprod \psi)(i)
  \defeq
  \begin{cases}
    \varphi(i), & 1\leq i\leq n,\\
    n'+\psi(i-n), & n<i\leq n+m.
  \end{cases}
\]

The monoidal unit is the unique typed interface
\[
  ([0],!),
  \qquad
  !:[0]\longrightarrow\operatorname{At}^{\Lambda}_{\Sigma}.
\]
The symmetry
\[
  \sigma_{U,V}:U \coprod V\longrightarrow V \coprod U
\]
is the type-preserving block transposition whose underlying function is
\[
  \sigma_{n,m}(i)
  \defeq
  \begin{cases}
    m+i, & 1\leq i\leq n,\\
    i-n, & n<i\leq n+m.
  \end{cases}
\]
\end{definition}

\begin{example}
Let $a,b,c,d\in\Sigma_0$ and let $U=([3],\tau)$ with
\[
  \tau(1)=a\multimap(b\otimes c),
  \qquad \tau(2)=a,
  \qquad \tau(3)=d.
\]
Then
\[
  \operatorname{ty}_{U}([3])
  =
  \langle a\multimap(b\otimes c),a,d\rangle,
  \qquad
  \operatorname{ty}_{U}(\{1,3\})
  =
  \langle a\multimap(b\otimes c),d\rangle.
\]
The order is inherited from $[3]$.
\end{example}

\begin{definition}
\label{def:lambda-discrete-interface-functor}
Define a functor
\[
  D_\Lambda:
  \mathbb F^\Lambda_\Sigma
  \longrightarrow
  \catname{CEHyp}_{\mathrm{s}}(\Sigma)
\]
as follows. For a typed interface
\[
  U=([n],\tau),
\]
let \(D_\Lambda(U)\) be the pointed discrete closed rooted
e-hypergraph with
\[
  V_{D_\Lambda(U)}
  =
  \{v_i\mid i\in[n]\},
  \qquad
  E_{D_\Lambda(U)}
  =
  \{\top_U\}.
\]

For a morphism
\[
  \varphi:U=([n],\tau)
  \longrightarrow
  V=([m],\upsilon)
\]
in \(\mathbb F^\Lambda_\Sigma\), define
\[
  D_\Lambda(\varphi):
  D_\Lambda(U)\longrightarrow D_\Lambda(V)
\]
by
\[
  D_\Lambda(\varphi)_V(v_i)
  =
  v_{\varphi(i)},
  \qquad
  D_\Lambda(\varphi)_E(\top_U)
  =
  \top_V.
\]
This is a strict homomorphism: preservation of vertex labels follows
from
\[
  \tau(i)=\upsilon(\varphi(i)),
\]
and the immediate-parent relation is preserved because
\[
  \top_V<^\mu v_{\varphi(i)}.
\]
These assignments preserve identities and composition and hence define
the stated functor. In particular,
\[
  D_\Lambda(0)=D_\Lambda([0],!)=\mathfrak T.
\]
\end{definition}

\begin{proposition}[Pushouts of closed rooted e-hypergraphs]
\label{prop:lambda-carrier-pushouts}
The pushout construction of \cref{prop:pushout_exists} restricts to
closed rooted e-hypergraphs.

Explicitly, let
\[
  X\xleftarrow{f}Z\xrightarrow{g}Y
\]
be a span in $\catname{CEHyp}_{\leq}(\Sigma)$ such that $Z$ is pointed
discrete, $f$ is strict and $g$ is anchored. Then its pushout exists in
both
\[
  \catname{CEHyp}_{\leq}(\Sigma)
  \qquad\text{and}\qquad
  \catname{CEHyp}_{\bullet}(\Sigma).
\]
If $g$ is strict, it is also a pushout in
$\catname{CEHyp}_{\mathrm{s}}(\Sigma)$.
\end{proposition}

\begin{proof}
Apply the construction of \cref{prop:pushout_exists}. Vertex
identifications preserve $l_V$ because both legs preserve vertex
types, and edge identifications preserve $l_E$ because both legs
preserve edge labels. The only edge of the pointed discrete apex is
its component root, so the construction neither identifies
$\boxsort$- or $\lambdasort$-edges nor changes their component
children. Consequently, the alternation and box-arity conditions are
preserved. The remainder of the proof is unchanged.
\end{proof}


Consequently,
\(\catname{CEHyp}_{\mathrm{s}}(\Sigma)\) is symmetric monoidal under
fibered coproduct, with unit \(\mathfrak T\). We fix the following
representatives for fibered coproducts of discrete interfaces. Given
\[
  U=([n],\tau)
  \qquad\text{and}\qquad
  V=([m],\upsilon),
\]
we take
\[
  D_\Lambda(U)\coprod_{\mathfrak T}D_\Lambda(V)
\]
to be the pointed discrete closed e-hypergraph with root \(\top\),
vertex set
\[
  \{v_1,\ldots,v_{n+m}\},
\]
and vertex labelling
\[
  l_V(v_i)=
  \begin{cases}
    \tau(i), & 1\leq i\leq n,\\
    \upsilon(i-n), & n<i\leq n+m.
  \end{cases}
\]
The coprojections send \(v_i\) in \(D_\Lambda(U)\) to \(v_i\), and
\(v_j\) in \(D_\Lambda(V)\) to \(v_{n+j}\); both send the root to
\(\top\). Thus, with these representatives,
\[
  D_\Lambda(U)\coprod_{\mathfrak T}D_\Lambda(V)
  =
  D_\Lambda(U\coprod V),
  \qquad
  \mathfrak T=D_\Lambda([0],!).
\]
Moreover, the fibered coproduct of
\(D_\Lambda(\varphi)\) and \(D_\Lambda(\psi)\) is
\(D_\Lambda(\varphi\coprod\psi)\), and the symmetry is the image under
\(D_\Lambda\) of the block transposition \(\sigma_{U,V}\).
Consequently, \(D_\Lambda\) is strict symmetric monoidal.

\begin{remark}
\label{rem:lambda-suppress-D}
The functor
\[
  D_\Lambda:
  \mathbb F^\Lambda_\Sigma
  \longrightarrow
  \catname{CEHyp}_{\mathrm{s}}(\Sigma)
\]
is fully faithful. Henceforth, we identify a typed interface
\[
  U=([n],\tau)
\]
with the pointed discrete closed e-hypergraph \(D_\Lambda(U)\), and
suppress \(D_\Lambda\) from the notation. We similarly write a morphism
\(\varphi\) of typed interfaces for its image \(D_\Lambda(\varphi)\).
That is, we will abbreviate $D_{\Lambda}([n], \tau)$ by writing $n$, so that the latter
includes both an ordered finite set of interface positions
and their atomic-type labelling.  
We write \(|n|\) for the number of positions of \(n\).
\end{remark}

\begin{definition}[Monogamous acyclic closed extended cospan]
\label{def:lambda-extended-cospan}
A \emph{closed extended cospan} from a typed interface \(n\) to a
typed interface \(k\) is a diagram
\[
  n
  \xrightarrow{f_{\mathrm{ext}}}
  n'
  \xrightarrow{f_{\mathrm{int}}}
  \mathcal G
  \xleftarrow{g_{\mathrm{int}}}
  k'
  \xleftarrow{g_{\mathrm{ext}}}
  k
\]
in \(\catname{CEHyp}_{\leq}(\Sigma)\), where
\(n,n',k',k\) are objects of \(\mathbb F^\Lambda_\Sigma\), regarded as
pointed discrete closed rooted e-hypergraphs according to
\cref{rem:lambda-suppress-D}, and \(\mathcal G\) is a closed rooted
e-hypergraph, subject to the following conditions:

\begin{enumerate}
  \item
  \(f_{\mathrm{ext}}:n\to n'\) and
  \(g_{\mathrm{ext}}:k\to k'\) are monomorphisms of typed interfaces;

  \item
  the composites
  \[
    f_{\mathrm{ext}}\fatsemi f_{\mathrm{int}}
    \qquad\text{and}\qquad
    g_{\mathrm{ext}}\fatsemi g_{\mathrm{int}}
  \]
  are strict homomorphisms;

  \item
  every strictly internal interface vertex, that is, every vertex of
  \(n'\) or \(k'\) outside the image of \(f_{\mathrm{ext}}\), respectively
  \(g_{\mathrm{ext}}\), is mapped below a non-root component of
  \(\mathcal G\).
  That is, for any strictly internal interface vertex \(x\), we have
  \[
    <^{\mu}_{\mathcal{G}}(f_{\mathrm{int},V}(x)) \neq \top_{\mathcal{G}} \qquad \text{and} \qquad <^{\mu}_{\mathcal{G}}(g_{\mathrm{int},V}(x)) \neq \top_{\mathcal{G}}
  \]
\end{enumerate}

A closed extended cospan is \emph{monogamous and acyclic}, or an
\emph{MDA closed extended cospan}, if moreover:

\begin{enumerate}
  \item the underlying directed hypergraph of \(\mathcal G\) is
  acyclic;

  \item \(f_{\mathrm{int}}\) and \(g_{\mathrm{int}}\) are
  monomorphisms;

  \item every vertex of \(\mathcal G\) has in-degree and out-degree at
  most one;

  \item a vertex of \(\mathcal G\) has in-degree zero if and only if it
  belongs to
  \(\operatorname{im}(f_{\mathrm{int},V})\);

  \item a vertex of \(\mathcal G\) has out-degree zero if and only if it
  belongs to
  \(\operatorname{im}(g_{\mathrm{int},V})\).
\end{enumerate}
We take the notion of isomorphism of closed extended cospans to be the same as in~\cref{def:isomorphic_ecospans}.
\end{definition}

\begin{definition}[Well-typed closed extended cospan]
\label{def:lambda-well-typed-extended-cospan}
Let
\[
  n
  \xrightarrow{f_{\mathrm{ext}}}
  n'
  \xrightarrow{f_{\mathrm{int}}}
  \mathcal G
  \xleftarrow{g_{\mathrm{int}}}
  k'
  \xleftarrow{g_{\mathrm{ext}}}
  k
\]
be an MDA closed extended cospan. Write
\[
  n'=([p],\tau),
  \qquad
  k'=([q],\upsilon),
\]
and let \(x_i\) and \(y_j\) denote the vertices corresponding,
respectively, to positions \(i\in[p]\) of \(n'\) and \(j\in[q]\) of
\(k'\).

Define
\[
\begin{aligned}
  \pi_f:[p]
  &\longrightarrow E_{\mathcal G,\cmpsortsubscript},
  &
  \pi_f(i)
  &\defeq
  <^{\mu}_{\mathcal G}
  \bigl(f_{\mathrm{int},V}(x_i)\bigr),
  \\
  \pi_g:[q]
  &\longrightarrow E_{\mathcal G,\cmpsortsubscript},
  &
  \pi_g(j)
  &\defeq
  <^{\mu}_{\mathcal G}
  \bigl(g_{\mathrm{int},V}(y_j)\bigr).
\end{aligned}
\]
These maps are well defined by the alternation condition. For a
$\cmpsort$-edge \(c\), the ordered subsets
\(\pi_f^{-1}(c)\subseteq[p]\) and
\(\pi_g^{-1}(c)\subseteq[q]\) record the input and output interface
positions belonging to \(c\).

A \(\boxsort\)-edge \(e\) is \emph{well typed} if, for every component
edge \(c\) that is an immediate child of \(e\),
\[
  \operatorname{ty}_{n'}\bigl(\pi_f^{-1}(c)\bigr)
  =
  l_V^*\bigl(s(e)\bigr),
  \qquad
  \operatorname{ty}_{k'}\bigl(\pi_g^{-1}(c)\bigr)
  =
  l_V^*\bigl(t(e)\bigr).
\]
Thus every component of \(e\) has the same ordered input and output
type words as \(e\) itself.

Let \(e\) be a \(\lambdasort\)-edge and let \(c\) be its unique
component child. We call \(e\) \emph{well-typed} if there exist
\[
A,B\in\operatorname{Ob}^{\Lambda}_{\Sigma} \qquad \text{and} \qquad X \in \operatorname{At}^{\Lambda}_{\Sigma}
\] 
such that
\[
  l_V^*\bigl(s(e)\bigr)
  =
  \operatorname{wd}(A),
  \qquad
  l_V^*\bigl(t(e)\bigr)
  =
  \operatorname{wd}(X\multimap B),
\]
and
\[
  \operatorname{ty}_{n'}\bigl(\pi_f^{-1}(c)\bigr)
  =
  \operatorname{wd}(A\otimes X),
  \qquad
  \operatorname{ty}_{k'}\bigl(\pi_g^{-1}(c)\bigr)
  =
  \operatorname{wd}(B).
\]
In other words, \(e\) represents the abstraction of a component with
type
\[
  A \otimes X \longrightarrow B
\]
to a morphism with type
\[
  A \longrightarrow X \multimap B.
\]
Since \(X\) is atomic,
\(\operatorname{wd}(X)=\langle X\rangle\); hence the input interface
of the component has exactly one additional position relative to the
source of \(e\).

The MDA closed extended cospan is \emph{well-typed} if every
\(\boxsort\)-edge and every \(\lambdasort\)-edge of \(\mathcal G\) is
well-typed.
\end{definition}


\begin{example}
Consider the extended cospan
\[
n \xrightarrow{f_{\mathrm{ext}}} n' \xrightarrow{f_{\mathrm{int}}} \mathcal G \xleftarrow{g_{\mathrm{int}}} k' \xleftarrow{g_{\mathrm{ext}}} k
\]
depicted below, where $A,B,C,X,Y$ are atomic types.
\[
\adjustbox{scale=0.65}{
  \tikzfig{./figures/lambda/closed_e_hypergraph_cospan_example}~.
}
\]
We note that this is a well-typed MDA closed extended cospan.
First, the degrees of vertices and the embeddings of the interfaces satisfy the MDA conditions and the carrier $\mathcal{G}$ is acyclic.
Then, note that
\[
\operatorname{ty}_{n'}(\pi_f^{-1}(E_2)) = \langle A,B,C \rangle, \qquad \operatorname{ty}_{n'}(\pi_f^{-1}(E_3)) = \langle A,B,C \rangle
\]
and
\[
\operatorname{ty}_{k'}(\pi_g^{-1}(E_2)) = \langle A,B,Y \rangle, \qquad \operatorname{ty}_{k'}(\pi_g^{-1}(E_3)) = \langle A,B,Y \rangle,
\]
therefore the $\boxsort$-edge $E_1$ is well-typed since
\[
l^*_{V}(s(E_1)) = \langle A,B,C \rangle, \qquad l^*_V(t(E_1)) = \langle A,B,Y \rangle.
\]
For the $\lambdasort$-edge $E_4$, we have
\[
l^{*}_V(s(E_4)) = \langle A, B \rangle, \qquad l^{*}_{V}(t(E_4)) = \langle X \multimap (A \otimes B \otimes Y) \rangle,
\]
and
\[
\operatorname{ty}_{n'}(\pi_f^{-1}(E_5)) = \langle A, B, X \rangle, \qquad \operatorname{ty}_{k'}(\pi_g^{-1}(E_5)) = \langle A, B, Y \rangle.
\]
To avoid a clash with these type names, write $X',A',B'$ for the objects required by \cref{def:lambda-well-typed-extended-cospan}. 
We take $X' \defeq X \in \operatorname{At}^{\Lambda}_{\Sigma}$, $A' \defeq A \otimes B \in \operatorname{Ob}^{\Lambda}_{\Sigma}$, and $B' \defeq A \otimes B \otimes Y \in \operatorname{Ob}^{\Lambda}_{\Sigma}$.
Hence, the $\lambdasort$-edge $E_4$ is well-typed.
Therefore, the extended cospan is well-typed, since all $\boxsort$- and $\lambdasort$-edges in its carrier are well-typed.
\end{example}

Henceforth, we will only consider well-typed MDA closed extended cospans, so we will refer to them simply as \emph{extended cospans}.

\begin{definition}[Category of closed extended cospans]
\label{def:lambda-extended-cospan-category}
The category
\[
  \catname{CMEHypI}_{\bullet}(\Sigma)
\]
has the objects of \(\mathbb F^\Lambda_\Sigma\) as its objects. A
morphism \(n\to k\) is an isomorphism class of well-typed MDA closed
extended cospans from \(n\) to \(k\).

Composition and identities  of extended cospans are
defined as in
\cref{def:rooted_ehypi}, using the pushouts of
\cref{prop:lambda-carrier-pushouts}. 
Composition preserves monogamy, acyclicity and it also preserves well-typedness: composition
does not change the source or target word of a hierarchical edge, nor
the ordered interface words assigned to any of its component children.

Analogously, the monoidal structure on objects is inherited from that of $\mathbb{F}^{\Lambda}_{\Sigma}$ (\textit{cf.}~\cref{def:rooted_ehypi}), and the tensor of morphisms is given by componentwise fibered coproducts of extended cospans.
Since the latter are taken up to isomorphism, the tensor is strict.
Finally, the symmetry is induced by the symmetry of $\coprod$ in $\mathbb{F}^{\Lambda}_{\Sigma}$, and therefore
\(\catname{CMEHypI}_{\bullet}(\Sigma)\) is a strict symmetric monoidal
category.
\end{definition}

\begin{definition}[Currying of extended cospans]
\label{def:lambda-currying-cospan}
Let
\[
  \mathbf G
  =
  \left(
  n
  \xrightarrow{f_{\mathrm{ext}}}
  n'
  \xrightarrow{f_{\mathrm{int}}}
  \mathcal G
  \xleftarrow{g_{\mathrm{int}}}
  k'
  \xleftarrow{g_{\mathrm{ext}}}
  k
  \right)
\]
be a well-typed MDA closed extended cospan. Write
\[
  n=([r],\tau),
  \qquad
  k=([q],\upsilon),
\]
and suppose that \(r\geq1\).

Define the typed prefix interface
\[
  \bar{n}\defeq([r-1],\bar{\tau}),
  \qquad
  \bar{\tau}(i)\defeq\tau(i)
  \quad (i\in[r-1]),
\]
and put
\[
  X
  \defeq
  \left\lVert
    \operatorname{ty}_{\bar{n}}([r-1])
  \right\rVert_\otimes,
  \qquad
  A\defeq\tau(r),
  \qquad
  B
  \defeq
  \left\lVert
    \operatorname{ty}_{k}([q])
  \right\rVert_\otimes.
\]
Thus
\[
  \operatorname{ty}_{n}([r])
  =
  \operatorname{wd}(A \otimes X),
  \qquad
  \operatorname{ty}_{k}([q])
  =
  \operatorname{wd}(B).
\]

Let
\[
  \mathbf 1_{X\multimap B}
  \defeq
  ([1],\chi_{X,B}),
  \qquad
  \chi_{X,B}(1)\defeq X\multimap B,
\]
be the singleton typed interface labelled by \(X \multimap B\).

Using the interface-block convention for extended cospans, choose a
representative of \(\mathbf G\) in which the external interfaces occur
as ordered blocks of \(n'\) and \(k'\).

Construct a closed rooted e-hypergraph
\(\mathcal G^\Lambda\) from \(\mathcal G\) by adjoining:

\begin{enumerate}
  \item a fresh component edge \(r_\Lambda\), which becomes the root of
  \(\mathcal G^\Lambda\);

  \item a fresh \(\lambdasort\)-edge \(e_\Lambda\);

  \item fresh vertices \(u_1,\ldots,u_{r-1}\) with
  \[
    l_V(u_i)=\tau(i)
    \qquad (i\in[r-1]);
  \]

  \item a fresh vertex \(z\) with
  \[
    l_V(z)=X\multimap B.
  \]
\end{enumerate}

The new edges are labelled and incident as follows:
\[
  l_E(r_\Lambda)=\cmpsort,
  \qquad
  s(r_\Lambda)=t(r_\Lambda)=\varepsilon,
\]
and
\[
  l_E(e_\Lambda)=\lambdasort,
  \qquad
  s(e_\Lambda)=\langle u_1,\ldots,u_{r-1}\rangle,
  \qquad
  t(e_\Lambda)=\langle z\rangle.
\]
The new immediate-parent relations are
\[
  r_\Lambda<^\mu e_\Lambda,
  \qquad
  r_\Lambda<^\mu u_i
  \quad(i\in[r-1]),
  \qquad
  r_\Lambda<^\mu z,
  \qquad
  e_\Lambda<^\mu\top_{\mathcal G}.
\]
All other incidence, labelling and immediate-parent data are inherited
from \(\mathcal G\). In particular, \(\top_{\mathcal G}\) is the unique
component child of \(e_\Lambda\).

Define
\[
  \eta_{\mathrm{in}}:
  \bar{n}\longrightarrow\mathcal G^\Lambda
\]
by sending the \(i\)-th vertex of \(\bar{n}\) to \(u_i\) and its root to
\(r_\Lambda\). Define
\[
  \eta_{\mathrm{out}}:
  \mathbf 1_{X\multimap B}
  \longrightarrow
  \mathcal G^\Lambda
\]
by sending its unique vertex to \(z\) and its root to \(r_\Lambda\).

Re-anchor the old internal-interface maps at the new root by putting
\[
  \widehat f_{\mathrm{int}}
  \defeq
  \left\langle
    r_\Lambda;(f_{\mathrm{int}})_V
  \right\rangle
  :
  n'\longrightarrow\mathcal G^\Lambda
\]
and
\[
  \widehat g_{\mathrm{int}}
  \defeq
  \left\langle
    r_\Lambda;(g_{\mathrm{int}})_V
  \right\rangle
  :
  k'\longrightarrow\mathcal G^\Lambda,
\]
where the old vertex maps are regarded as maps into
\(\mathcal G^\Lambda\) through the inclusion
\(\mathcal G\hookrightarrow\mathcal G^\Lambda\).

The \emph{currying} of \(\mathbf G\) is the closed extended cospan
\[
  \Lambda(\mathbf G)
  \defeq
  \left(
  \bar{n}
  \xrightarrow{\iota_{\bar{n}}}
  \bar{n}\coprod n'
  \xrightarrow{
    [\eta_{\mathrm{in}},\widehat f_{\mathrm{int}}]
  }
  \mathcal G^\Lambda
  \xleftarrow{
    [\eta_{\mathrm{out}},\widehat g_{\mathrm{int}}]
  }
  \mathbf 1_{X\multimap B}\coprod k'
  \xleftarrow{\iota_{\mathbf 1_{X\multimap B}}}
  \mathbf 1_{X\multimap B}
  \right),
\]
where the outer maps are the first coprojections.

Thus,
\[
  \Lambda(\mathbf G):
  \bar{n}\longrightarrow\mathbf 1_{X\multimap B}
\]
represents the currying operation
\[
  A\otimes X\longrightarrow B
  \quad\longmapsto\quad
  A\longrightarrow X\multimap B.
\]
\end{definition}

\begin{example}
Below is an example of currying a closed extended cospan.
\[
\adjustbox{scale=0.65}{
  \tikzfig{./figures/lambda/lambda_example}~.
}
\]
\end{example}

\begin{remark}
The join of MDA closed extended cospans is defined in the same way as in~\cref{def:join_of_extended_cospans}.
\end{remark}


\begin{definition}[Concrete interpretation of closed-monoidal terms]
\label{def:lambda-concrete-interpretation}
For \(A\in\operatorname{Ob}^{\Lambda}_{\Sigma}\), let
\(\partial A\) denote the typed interface determined by
\(\operatorname{wd}(A)\). Explicitly, if
\[
  \operatorname{wd}(A)=\langle A_1,\ldots,A_r\rangle,
\]
then
\[
  \partial A
  \defeq
  ([r],i\mapsto A_i).
\]
In particular,
\[
  \partial I=0,
  \qquad
  \partial(A\otimes B)
  =
  \partial A\coprod\partial B,
\]
and \(\partial(A\multimap B)\) is the singleton interface labelled
\(A\multimap B\).

We define the interpretation of raw closed-monoidal terms recursively. A term
\[
  t:A\longrightarrow B
\]
is interpreted as a well-typed MDA closed extended cospan
\[
  \widehat t:
  \partial A\longrightarrow\partial B.
\]

Let \(f\in\Sigma_1\) be a generator with
\[
  \operatorname{dom}(f)
  =
  \langle A_1,\ldots,A_n\rangle,
  \qquad
  \operatorname{cod}(f)
  =
  \langle B_1,\ldots,B_k\rangle.
\]
Its interpretation is the one-edge closed extended cospan
\[
  \widehat f
  \defeq
  \adjustbox{scale=0.65}{
    \tikzfig{./figures/lambda/generator_interpretation}
  }.
\]

For \(A\in\operatorname{At}^{\Lambda}_{\Sigma}\) and
\(B\in\operatorname{Ob}^{\Lambda}_{\Sigma}\), the evaluation term
\[
  \ev_{A,B}:
  (A\multimap B)\otimes A\longrightarrow B
\]
is interpreted as
\[
  \widehat{\ev_{A,B}}
  \defeq
  \adjustbox{scale=0.65}{
    \tikzfig{./figures/lambda/eval_generator_interpretation}
  },
\]
where the ordered output vertices are labelled by
\(\operatorname{wd}(B)\).

Primitive identities and symmetries are interpreted by the corresponding
identities and symmetries in
\(\catname{CMEHypI}_{\bullet}(\Sigma)\):
\[
  \widehat{\id_I}=\id_{\partial I},
  \qquad
  \widehat{\id_T}=\id_{\partial T},
  \qquad
  \widehat{\sym_{T,U}}
  =
  \sym_{\partial T,\partial U}.
\]

The remaining constructors are interpreted recursively by
\[
  \widehat{f\fatsemi g}
  \defeq
  \widehat f\fatsemi\widehat g,
  \qquad
  \widehat{f\otimes g}
  \defeq
  \widehat f\otimes\widehat g,
\]
and
\[
  \widehat{\Lambda_{A,X,B}(f)}
  \defeq
  \Lambda(\widehat f).
\]
Here \(X\in\operatorname{At}^{\Lambda}_{\Sigma}\), so the final input
position of
\[
  \widehat h:
  \partial A\coprod\partial X \longrightarrow\partial B
\]
is a single vertex, and the currying construction of
\cref{def:lambda-currying-cospan} applies.

For terms with joins, define additionally
\[
  \widehat{f_1\join\cdots\join f_n}
  \defeq
  \widehat{f_1}\join\cdots\join\widehat{f_n}
  \qquad(n\geq2),
\]
where the right-hand side is the $n$-ary join of closed extended
cospans of \cref{def:join_of_extended_cospans}. 
This gives interpretation functions
\[
  \operatorname{Tm}^{\Lambda}_{\Sigma}(A,B)
  \longrightarrow
  \catname{CMEHypI}_{\bullet}(\Sigma)
  \bigl(\partial A,\partial B\bigr)
\]
and
\[
  \operatorname{Tm}^{\Lambda,\join}_{\Sigma}(A,B)
  \longrightarrow
  \catname{CMEHypI}_{\bullet}(\Sigma)
  \bigl(\partial A,\partial B\bigr).
\]
\end{definition}

\section{Convex rewriting}

The construction of extended pushout complements from
\cref{def:extended_complement} applies to closed rooted
e-hypergraphs. In this setting all interfaces and interface
isomorphisms are typed, all pushouts are taken in
\(\catname{CEHyp}_{\bullet}(\Sigma)\), and well-typedness means
\cref{def:lambda-well-typed-extended-cospan}. If \(S\) is a typed
subinterface of \(n'\), then \(n'\setminus S\) denotes the ordered
complement of \(S\), equipped with the restricted type labelling.
Again, under the convention as in~\cref{remark:interface_convention}, these complements might be regarded as omission of certain components of a fibered coproduct.
The descendant-closure condition applies equally to
\(\boxsort\)- and \(\lambdasort\)-edges.

\begin{definition}[Extended boundary pushout complement]
\label{def:lambda-boundary-complement}
Consider well-typed MDA closed extended cospans
\[
  i
  \xrightarrow{f_{\mathrm{ext}}^{i}}
  i'
  \xrightarrow{f_{\mathrm{int}}^{i}}
  \mathcal L
  \xleftarrow{f_{\mathrm{int}}^{j}}
  j'
  \xleftarrow{f_{\mathrm{ext}}^{j}}
  j
\]
and
\[
  n
  \xrightarrow{g_{\mathrm{ext}}^{n}}
  n'
  \xrightarrow{g_{\mathrm{int}}^{n}}
  \mathcal G
  \xleftarrow{g_{\mathrm{int}}^{k}}
  k'
  \xleftarrow{g_{\mathrm{ext}}^{k}}
  k,
\]
and let
\[
  m:\mathcal L\longrightarrow\mathcal G
\]
be an anchored homomorphism. Write
\[
  f
  \defeq
  \bigl[
    f_{\mathrm{ext}}^{i}\fatsemi f_{\mathrm{int}}^{i},
    f_{\mathrm{ext}}^{j}\fatsemi f_{\mathrm{int}}^{j}
  \bigr]
  :
  i\coprod j\longrightarrow\mathcal L.
\]

Suppose that
\[
  i \coprod j
  \xrightarrow{[c_1,c_2]}
  \mathcal L^\bot
  \xrightarrow{l^\bot}
  \mathcal G
\]
together with the retained interface maps
\[
  g_n^\bot:
  n'\setminus S_n\longrightarrow\mathcal L^\bot,
  \qquad
  g_k^\bot:
  k'\setminus S_k\longrightarrow\mathcal L^\bot
\]
and the typed block isomorphisms
\[
  \alpha_n:i'\setminus i\longrightarrow S_n,
  \qquad
  \alpha_k:j'\setminus j\longrightarrow S_k
\]
is an extended pushout complement of \((f,m)\) in the sense of
\cref{def:extended_complement}, interpreted in
\(\catname{CEHyp}_{\bullet}(\Sigma)\).

Define the re-anchored interface maps
\[
  \widehat c_1
  \defeq
  \left\langle
    \top_{\mathcal L^\bot};(c_1)_V
  \right\rangle
  :
  i\longrightarrow\mathcal L^\bot,
  \qquad
  \widehat c_2
  \defeq
  \left\langle
    \top_{\mathcal L^\bot};(c_2)_V
  \right\rangle
  :
  j\longrightarrow\mathcal L^\bot.
\]
Thus \(\widehat c_1\) and \(\widehat c_2\) have the same vertex maps as
\(c_1\) and \(c_2\), respectively, but send their roots to
\(\top_{\mathcal L^\bot}\).

The extended pushout complement is an \emph{extended boundary pushout
complement} if the following conditions hold:

\begin{enumerate}
  \item \(m\) is a monomorphism;

  \item \([c_1,c_2]\) is a monomorphism;

  \item if
  \[
    m_E(\top_{\mathcal L})=\top_{\mathcal G},
  \]
  then
  \[
    n \coprod j
    \xrightarrow{
      \overline g_{\mathrm{ext}}^{\,n} \coprod \id_j
    }
    (n'\setminus S_n) \coprod j
    \xrightarrow{
      [g_n^\bot,\widehat c_2]
    }
    \mathcal L^\bot
    \xleftarrow{
      [g_k^\bot,\widehat c_1]
    }
    (k'\setminus S_k) \coprod i
    \xleftarrow{
      \overline g_{\mathrm{ext}}^{\,k} \coprod \id_i
    }
    k \coprod i
  \]
  is a well-typed MDA closed extended cospan;

  \item if
  \[
    m_E(\top_{\mathcal L})\neq\top_{\mathcal G},
  \]
  then
  \[
    n
    \xrightarrow{
      \overline g_{\mathrm{ext}}^{\,n}
      \fatsemi\iota_{n'\setminus S_n}
    }
    (n'\setminus S_n) \coprod j
    \xrightarrow{
      [g_n^\bot,\widehat c_2]
    }
    \mathcal L^\bot
    \xleftarrow{
      [g_k^\bot,\widehat c_1]
    }
    (k'\setminus S_k) \coprod i
    \xleftarrow{
      \overline g_{\mathrm{ext}}^{\,k}
      \fatsemi\iota_{k'\setminus S_k}
    }
    k
  \]
  is an MDA closed extended cospan, not necessarily well typed. Here
  \(\iota_{n'\setminus S_n}\) and
  \(\iota_{k'\setminus S_k}\) are the first coprojections into the
  displayed internal interfaces.
\end{enumerate}
\end{definition}

\begin{proposition}
\label{prop:lambda-boundary-complement-unique}
The extended boundary pushout complement of \cref{def:lambda-boundary-complement} is unique up to isomorphism.
\end{proposition}
\begin{proof}
This is analogous to~\cref{prop:extended_boundary_complement_unique}.
\end{proof}

\begin{definition}[Convex closed EDPOI rewriting]
\label{def:lambda-convex-edpoi-rewriting}
A \emph{closed rewrite rule} is a pair of well-typed MDA closed
extended cospans with common external interfaces. In folded form, it is
displayed as
\[
  \mathcal L
  \xleftarrow{f_{\mathrm{int}}}
  i'\coprod j'
  \xleftarrow{f_{\mathrm{ext}}}
  i\coprod j
  \xrightarrow{p_{\mathrm{ext}}}
  i''\coprod j''
  \xrightarrow{p_{\mathrm{int}}}
  \mathcal R.
\]
Put
\[
  f
  \defeq
  f_{\mathrm{ext}}\fatsemi f_{\mathrm{int}},
  \qquad
  p
  \defeq
  p_{\mathrm{ext}}\fatsemi p_{\mathrm{int}}.
\]

Let
\[
  n
  \longrightarrow
  n'
  \longrightarrow
  \mathcal G
  \longleftarrow
  k'
  \longleftarrow
  k
\]
be a well-typed MDA closed extended cospan. A \emph{convex closed
EDPOI rewrite} of \(\mathcal G\) by
\(\langle\mathcal L,\mathcal R\rangle\) consists of an EDPOI diagram
\[
\begin{tikzcd}
  \mathcal L
    \arrow[d,"m"']
  &
  i\coprod j
    \arrow[l,"f"']
    \arrow[r,"p"]
    \arrow[d,"c"]
  &
  \mathcal R
    \arrow[d,"\widetilde\iota_{\mathcal R}"]
  \\
  \mathcal G
  &
  \mathcal L^\bot
    \arrow[l,"l^\bot"]
    \arrow[r,"\widetilde\iota_{\bot}"']
  &
  \mathcal H
\end{tikzcd}
\]
in \(\catname{CEHyp}_{\bullet}(\Sigma)\), satisfying the following
conditions:

\begin{enumerate}
  \item both squares are pushouts;

  \item the image of the anchored match
  \(m:\mathcal L\to\mathcal G\) is a convex sub-e-hypergraph of
  \(\mathcal G\);

  \item the left-hand square, together with its retained and affected
  typed-interface data, is an extended boundary pushout complement in
  the sense of \cref{def:lambda-boundary-complement};

  \item the internal interfaces and their maps into \(\mathcal H\) are
  obtained from the EDPOI construction of \cref{def:dpoi_rooted}.
  Explicitly, if \(S_n\subseteq n'\) and \(S_k\subseteq k'\) are the
  affected subinterfaces, then the internal interfaces of the result
  are
  \[
    n''
    \defeq
    (n'\setminus S_n)\coprod(i''\setminus i),
    \qquad
    k''
    \defeq
    (k'\setminus S_k)\coprod(j''\setminus j).
  \]
  The retained interface maps are composed with
  \(\widetilde\iota_{\bot}\), while the strictly internal interface
  maps of \(\mathcal R\) are composed with
  \(\widetilde\iota_{\mathcal R}\) and re-anchored at
  \(\top_{\mathcal H}\).
\end{enumerate}

The resulting well-typed MDA closed extended cospan is
\[
  n
  \longrightarrow
  n''
  \longrightarrow
  \mathcal H
  \longleftarrow
  k''
  \longleftarrow
  k.
\]
We write
\[
  \mathcal G
  \mathrel{
    \widetilde{\rightsquigarrow}_{
      \langle\mathcal L,\mathcal R\rangle
    }
  }
  \mathcal H
\]
to denote that $\mathcal{G}$ rewrites to $\mathcal{H}$ via the convex closed EDPOI rewriting.
\end{definition}

\begin{definition}[Convex match]
\label{def:lambda-convex-match}
Let $\mathcal L$ and $\mathcal G$ be well-typed MDA closed extended
cospans. A \emph{convex match} of $\mathcal L$ in $\mathcal G$ is an
anchored monomorphism
\[
  m:\mathcal L\longrightarrow\mathcal G
\]
whose image is a convex sub-e-hypergraph of $\mathcal G$, together with
a choice of extended boundary pushout complement of $m$ in the sense of
\cref{def:lambda-boundary-complement}. These are exactly the match data
of a convex closed EDPOI rewrite
(\cref{def:lambda-convex-edpoi-rewriting}).

Recall that condition~(2) of \cref{def:extended_complement} is part of
the complement data. Hence, the image of a convex match is closed under
descendants below every matched non-root edge, and a convex match
satisfies the hypotheses of \cref{lemma:lambda-convex-decomposition}.
\end{definition}

The following lemma is the closed analogue of~\cref{lemma:extended_convex_decomposition}.
This is the main workhorse for the proof of the equivalence of rewriting on terms and EDPOI rewriting on closed e-hypergraphs and the payoff of convex matches.
While the structure of a term determines the structure of the corresponding extended cospan because the interpretation is compositional, convex decomposition determines the structure of the extended cospan given a convex match.

\begin{lemma}[Closed convex decomposition]
\label{lemma:lambda-convex-decomposition}
Let \(\mathbf L:i\to j\) and \(\mathbf F\) be well-typed MDA closed
extended cospans, with carriers \(\mathcal L\) and \(\mathcal F\), and let
\[
  m:\mathcal L\longrightarrow\mathcal F
\]
be an anchored monomorphism whose image is convex in the underlying
directed hypergraph of \(\mathcal F\). Suppose that the image of \(m\)
is closed under descendants below every matched non-root edge, as in
\cref{lemma:extended_convex_decomposition}, with the condition applied
to both \(\boxsort\)- and \(\lambdasort\)-edges.

Then the recursive decomposition of
\cref{lemma:extended_convex_decomposition} extends to the closed
setting. More precisely, there are a well-typed MDA closed extended
cospan \(\mathbf D\), with carrier \(\mathcal D\), constructed from
\(\mathbf L\) and other well-typed MDA closed extended cospans by
composition, tensor product, join and abstraction, and an isomorphism
\[
  \mathbf{\Phi}:\mathbf D\xrightarrow{\cong}\mathbf F
\]
such that the canonical occurrence
\[
  \iota_{\mathcal L}:
  \mathcal L \longrightarrow \mathcal{D}
\]
of \(\mathcal L\) in \(\mathcal D\) satisfies
\[
  \iota_{\mathcal L}\fatsemi \phi = m,
  \tag{\(\ast\)}
  \label{eq:lambda-convex-decomposition-occurrence}
\]
where \(\phi\) is the carrier component of \(\mathbf{\Phi}\).

The root and \(\boxsort\)-parent clauses are those of
\cref{lemma:extended_convex_decomposition}. The only additional
recursive clause is the following. Suppose that
\[
  a\defeq m_E(\top_{\mathcal L})\neq\top_{\mathcal F}
\]
and that the unique immediate parent \(p\) of \(a\) is a
\(\lambdasort\)-edge. 
Then the anchor \(a\) is the unique component child of \(p\).
It determines a well-typed component cospan
\[
  \mathbf F_a:\partial(X\otimes A)\longrightarrow\partial B,
\]
where
\[
  X,B\in\operatorname{Ob}^{\Lambda}_{\Sigma},
  \qquad
  A\in\operatorname{At}^{\Lambda}_{\Sigma},
\]
and \(m\) corestricts to a convex, descendant-closed anchored
monomorphism
\[
  m_a:\mathcal L\longrightarrow\mathcal F_a.
\]

Applying the root clause inside \(\mathbf F_a\) gives a typed interface
\(q\), well-typed MDA closed extended cospans
\[
  \mathbf C_1^a:
  \partial(X\otimes A)\longrightarrow q\coprod i,
  \qquad
  \mathbf C_2^a:
  q\coprod j\longrightarrow \partial B,
\]
and an isomorphism
\[
  \mathbf{\Phi}_{a}:
  \mathbf D_a
  \defeq
  \mathbf C_1^a
  \fatsemi
  \bigl(\mathbf 1_q\otimes\mathbf L\bigr)
  \fatsemi
  \mathbf C_2^a
  \xrightarrow{\cong}
  \mathbf F_a
\]
such that
\[
  \iota_{\mathcal L}^a
  \fatsemi
  \phi_a
  =
  m_a,
  \tag{\(\dagger\)}
  \label{eq:lambda-component-occurrence}
\]
where \(\iota^a_{\mathcal L}\) is the canonical occurrence of
\(\mathcal L\) in the carrier of \(\mathbf D_a\).

Consequently, after identifying the subcospan generated by \(p\) with
\(\Lambda(\mathbf F_a)\), abstraction induces an isomorphism
\[
  \mathbf{\Phi}_p:
  \Lambda(\mathbf D_a)
  \xrightarrow{\cong}
  \Lambda(\mathbf F_a)
\]
which carries the canonical occurrence of \(\mathcal L\) to the
corestriction \(m_p\) of \(m\):
\[
  \iota_{\mathcal L}^p
  \fatsemi
  \phi_p
  =
  m_p.
  % \tag{\(\ddagger\)}
  % \label{eq:lambda-parent-occurrence}
\]

If \(p\) is itself nested, the same construction is applied recursively
to the occurrence of the subcospan generated by \(p\) in its enclosing
component. Each application moves one level towards
\(\top_{\mathcal F}\), so the process terminates.
\end{lemma}

\begin{proof}
The first two clauses are those of \cref{lemma:extended_convex_decomposition}, so it
remains to consider a \(\lambdasort\)-parent \(p\). Let
\[
  \jmath_{\mathbf D_a}:
  \mathcal D_a
  \longrightarrow
  \Lambda(\mathcal D_a)
  \quad\text{and}\quad
  \jmath_{\mathbf F_a}:
  \mathcal F_a
  \longrightarrow
  \Lambda(\mathcal F_a)
\]
be the canonical inclusions of the component carriers into the carriers
\(\Lambda(\mathcal D_a)\) and \(\Lambda(\mathcal F_a)\) of
\(\Lambda(\mathbf D_a)\) and \(\Lambda(\mathbf F_a)\)
(\cref{def:lambda-currying-cospan}).
The abstraction construction applies \(\phi_a\) to the component and
acts identically on the newly introduced \(\lambdasort\)-edge and its
ports. Consequently,
\begin{equation}
  \jmath_{\mathbf D_a}
  \fatsemi
  \Lambda(\phi_a)
  =
  \phi_a
  \fatsemi
  \jmath_{\mathbf F_a}.
  \tag{\(\star\)}
  \label{eq:lambda-convex-decomposition-abstraction-commutes}
\end{equation}
The canonical occurrence and the corestricted match are respectively
\[
  \iota_{\mathcal L}^p
  =
  \iota_{\mathcal L}^a
  \fatsemi
  \jmath_{\mathbf D_a},
  \qquad
  m_p
  =
  m_a
  \fatsemi
  \jmath_{\mathbf F_a}.
\]
Using \eqref{eq:lambda-component-occurrence} and \eqref{eq:lambda-convex-decomposition-abstraction-commutes}, we obtain
\[
\begin{aligned}
  \iota_{\mathcal L}^p
  \fatsemi
  \Lambda(\phi_a)
  &=
  \iota_{\mathcal L}^a
  \fatsemi
  \jmath_{\mathbf D_a}
  \fatsemi
  \Lambda(\phi_a)
  \\
  &=
  \iota_{\mathcal L}^a
  \fatsemi
  \phi_a
  \fatsemi
  \jmath_{\mathbf F_a}
  \\
  &=
  m_a\fatsemi\jmath_{\mathbf F_a}
  \\
  &=
  m_p.
\end{aligned}
\]
Composing, if necessary, with the canonical 
isomorphism from \(\Lambda(\mathbf F_a)\) to the subcospan generated by
\(p\) gives \(\mathbf{\Phi}_p\) and preserves this equation.

Repeating the appropriate \(\boxsort\)- or \(\lambdasort\)-parent
clause moves the decomposition outwards through the finite hierarchy.
Composition, tensor product, join and abstraction preserve the
displayed occurrence equation. The resulting isomorphism
\(\mathbf{\Phi}:\mathbf D\cong\mathbf F\) therefore satisfies
\[
  \iota_{\mathcal L}\fatsemi\phi = m,
\]
as required.
\end{proof}

Again, the intuition is that the decomposition inside hierarchical edges propagates upwards and moreover that match factors through the decomposition.

\section{Structural adequacy}
\label{sec:lambda-structural-adequacy}
We now show that the interpretation identifies exactly the intended
structural equations. 
First, we show that the join-free interpretation is faithful modulo the symmetric monoidal
equations. 
Then, we represent a term with joins by the finite set of its
join-free summands and represent an extended cospan by the set of its
complete resolutions. Comparing these two normal forms reduces the
general adequacy result to the join-free case.

\begin{proposition}[Faithfulness of the join-free interpretation]
\label{prop:lambda-join-free-faithfulness}
For raw closed-monoidal terms
\[
  s,t\in\operatorname{Tm}^{\Lambda}_{\Sigma}(A,B),
\]
one has
\[
  s\equiv_0 t
  \quad\Longleftrightarrow\quad
  \widehat s \cong \widehat t.
\]
\end{proposition}

\begin{proof}
The forward implication follows by induction on the derivation of
\(s\equiv_0t\). The interpretation takes the symmetric monoidal
constructors to those of the strict symmetric monoidal category
\(\catname{CMEHypI}_{\bullet}(\Sigma)\), so it validates the symmetric
monoidal equations up to isomorphism. Moreover, composition, tensor
product and abstraction preserve isomorphisms of closed extended
cospans. Hence
\[
  \widehat s\cong\widehat t.
\]

For the converse, suppose that
\[
  \phi:\widehat s\overset{\cong}{\longrightarrow}\widehat t
\]
is an isomorphism of closed extended cospans. Since \(s\) and \(t\) are
join-free, their interpretations contain no \(\boxsort\)-edges. We
proceed by induction on the maximum nesting depth of abstractions in
\(s\) and \(t\). By the recursive definition of the interpretation,
this is precisely the maximum nesting depth of \(\lambdasort\)-edges.
Since \(\phi\) preserves edge labels and immediate parents, \(s\) and
\(t\) have the same maximum depth.

If this depth is zero, neither interpretation contains a
\(\lambdasort\)-edge. Removing the root from each carrier therefore
gives isomorphic ordinary typed MDA hypergraph cospans over the
ordinary monoidal signature \(\Gamma_\Sigma\). These are precisely the
ordinary hypergraph interpretations of \(s\) and \(t\). Faithfulness of
the ordinary symmetric monoidal interpretation,
\cref{prop:faithful_interpretation}, applied in the multi-sorted setting
with sorts \(\operatorname{At}^{\Lambda}_{\Sigma}\), gives%
\footnote{Strictly speaking, faithfulness applies to the join-free and
abstraction-free fragment of
\(\operatorname{Tm}^{\Lambda}_{\Sigma}/{\equiv_0}\), regarded as
\(\mathbf S_{\Gamma_\Sigma}\).}
\[
  s\equiv_0t.
\]

Now suppose that the result holds below abstraction depth \(d>0\), and
that \(s\) and \(t\) have maximum abstraction depth \(d\). Consider the
outermost \(\lambdasort\)-edges of \(\widehat s\), namely those whose
immediate parent is \(\top_{\widehat s}\). For every such edge \(e\),
introduce a fresh generator \(x_e\) having the same ordered input and
output types as \(e\). Let
\[
  \Gamma_\phi
  \defeq
  \Gamma_\Sigma
  \sqcup
  \left\{
    x_e
    \;\middle|\;
    e\text{ is an outermost \(\lambdasort\)-edge of }\widehat s
  \right\}.
\]

For every outermost \(\lambdasort\)-edge \(e\) of \(\widehat s\), remove
its unique component child and all its descendants, discard the
corresponding strictly internal interface blocks, and regard \(e\) as
an ordinary edge labelled by \(x_e\). Perform the same construction in
\(\widehat t\) at the edge \(\phi_E(e)\), using the same generator
\(x_e\). Finally, remove the roots. This produces ordinary typed MDA
hypergraph cospans over \(\Gamma_\phi\), and \(\phi\) restricts to an
isomorphism between them.

By the recursive definition of the interpretation, every outermost
\(\lambdasort\)-edge \(e\) of \(\widehat s\) corresponds to a unique
occurrence of an abstraction \(\Lambda(h_e)\) in \(s\). Similarly,
\(\phi_E(e)\) corresponds to a unique occurrence
\(\Lambda(h'_e)\) in \(t\). Replace these occurrences by \(x_e\) to
obtain the flattened terms \(\bar s\) and \(\bar t\). Their ordinary
hypergraph interpretations are precisely the two flattened cospans
constructed above. Faithfulness of the ordinary many-sorted
hypergraph interpretation (\emph{cf.}~\cref{prop:faithful_interpretation}) therefore gives
\[
  [\bar s]_{\mathsf{SMC}}
  =
  [\bar t]_{\mathsf{SMC}}
\]
in the free symmetric monoidal category \(\mathbf S_{\Gamma_\phi}\), where
\([-]_{\mathsf{SMC}}\) denotes the class of a raw term modulo the
symmetric monoidal equations.

For every outermost \(\lambdasort\)-edge \(e\), the isomorphism
\(\phi\) restricts to an isomorphism of closed extended cospans
\[
  \widehat{h_e}\cong\widehat{h'_e}.
\]
These bodies have abstraction depth strictly smaller than \(d\), so
the induction hypothesis gives
\[
  h_e\equiv_0h'_e.
\]
Since \(\equiv_0\) is a congruence, it follows that
\[
  \Lambda(h_e)\equiv_0\Lambda(h'_e).
\]

As above, \(\mathbf S_{\Gamma_\phi}\) is the free symmetric monoidal
category on the multi-sorted signature \(\Gamma_\phi\). Let
\(\operatorname{Tm}^{\Lambda}_{\Sigma}/{\equiv_0}\) be the symmetric
monoidal category whose objects are
\(\operatorname{Ob}^{\Lambda}_{\Sigma}\) and whose morphisms are
\[
  \bigl(
    \operatorname{Tm}^{\Lambda}_{\Sigma}/{\equiv_0}
  \bigr)(A,B)
  \defeq
  \operatorname{Tm}^{\Lambda}_{\Sigma}(A,B)/{\equiv_0}.
\]
Let \([-]_0\) denote the corresponding quotient map.

Send each colour \(T\in\operatorname{At}^{\Lambda}_{\Sigma}\) to the
object \(T\), and hence send a word of colours to its tensor
realisation. The assignments
\[
\begin{aligned}
  x_e &\longmapsto [\Lambda(h_e)]_0,\\
  f   &\longmapsto [f]_0
        &&\bigl(f\in(\Gamma_\Sigma)_1\bigr)
\end{aligned}
\]
extend, by freeness, to a strict symmetric monoidal functor
\[
  F_\phi:
  \mathbf S_{\Gamma_\phi}
  \longrightarrow
  \operatorname{Tm}^{\Lambda}_{\Sigma}/{\equiv_0}.
\]

Replacing every \(x_e\) in \(\bar s\) by \(\Lambda(h_e)\) recovers
\(s\), and hence
\[
  F_\phi([\bar s]_{\mathsf{SMC}})=[s]_0.
\]
Let \(t^\phi\) be the closed \(\Sigma\)-term obtained by replacing
each occurrence of \(x_e\) in \(\bar t\) by \(\Lambda(h_e)\). Then
\[
  F_\phi([\bar t]_{\mathsf{SMC}})=[t^\phi]_0.
\]
On the other hand, \(t\) is recovered from \(\bar t\) by replacing
each \(x_e\) by \(\Lambda(h'_e)\). Since
\[
  \Lambda(h_e)\equiv_0\Lambda(h'_e)
\]
for every \(e\), congruence gives
\[
  t^\phi\equiv_0t.
\]

Applying \(F_\phi\) to
\(
[\bar s]_{\mathsf{SMC}}=[\bar t]_{\mathsf{SMC}}
\)
therefore gives
\[
  [s]_0
  =
  F_\phi([\bar s]_{\mathsf{SMC}})
  =
  F_\phi([\bar t]_{\mathsf{SMC}})
  =
  [t^\phi]_0
  =
  [t]_0.
\]
Consequently,
\[
  s\equiv_0t.
\]
\end{proof}

% This section isolates the structural part of the correspondence between
% closed semilattice-enriched terms and closed e-hypergraphs.  The result
% below compares structural equality of raw terms with structural EDPOI
% equivalence of their concrete interpretations.  Together with the
% reification of \cref{lemma:lambda-reification}, proved in
% \cref{sec:lambda-rewriting-correspondence}, it yields the corresponding
% isomorphism of quotients.

% We use the structural congruences $\equiv_0$ and
% $\equiv_{\mathsf{str}}$ of
% \cref{def:lambda-structural-congruences}. They contain no closed-category
% equations: $\beta$, $\eta$ and currying naturality may instead be treated as
% operational rules.

\subsection{Normal forms of raw closed-monoidal terms}

For types $A,B$, write
\[
 J_{\Sigma}(A,B)
 \defeq
 \operatorname{Tm}^{\Lambda}_{\Sigma}(A,B)/{\equiv_{0}}
\]
for the hom-sets of the category
$\operatorname{Tm}^{\Lambda}_{\Sigma}/{\equiv_0}$ introduced in the proof of
\cref{prop:lambda-join-free-faithfulness}.
Let $\mathcal P^{+}_{\mathrm{fin}}(X)$ denote the set of finite non-empty
subsets of a set $X$.  We define
\[
 N_{A,B}:
 \operatorname{Tm}^{\Lambda,\join}_{\Sigma}(A,B)
 \longrightarrow
 \mathcal P^{+}_{\mathrm{fin}}(J_{\Sigma}(A,B))
\]
by structural recursion.  We omit the type subscripts and set
\begin{align*}
 N(g)
   &\defeq \{[g]_{0}\},\\
 N(f_1\join\cdots\join f_n)
   &\defeq N(f_1)\cup\cdots\cup N(f_n),\\
 N(s\fatsemi t)
   &\defeq
     \bigl\{[u\fatsemi v]_{0}
       \mathrel{|}[u]_{0}\in N(s),\ [v]_{0}\in N(t)\bigr\},\\
 N(s\otimes t)
   &\defeq
     \bigl\{[u\otimes v]_{0}
       \mathrel{|}[u]_{0}\in N(s),\ [v]_{0}\in N(t)\bigr\},\\
 N(\Lambda(s))
   &\defeq
     \bigl\{[\Lambda(u)]_{0}
       \mathrel{|}[u]_{0}\in N(s)\bigr\}.
\end{align*}
Here $g$ ranges over atomic terms: identities, symmetries, generators
and evaluations.  

\begin{lemma}[Term normal forms]
\label{lemma:lambda-term-normal-forms}
For raw closed-monoidal terms with joins
\[
  s,t\in
  \operatorname{Tm}^{\Lambda,\join}_{\Sigma}(A,B),
\]
one has
\[
  s\equiv_{\mathsf{str}}t
  \quad\Longleftrightarrow\quad
  N(s)=N(t).
\]
\end{lemma}

\begin{proof}
Every generating equation of \(\equiv_{\mathsf{str}}\) preserves \(N\).
For the symmetric monoidal equations, expand each operand into its
finite set of join-free classes. For every choice of join-free
summands, the corresponding terms on the two sides are
\(\equiv_0\)-equivalent.

For example, bifunctoriality gives
\[
  (s\otimes t)\fatsemi(u\otimes v)
  \equiv_{\mathsf{str}}
  (s\fatsemi u)\otimes(t\fatsemi v).
\]
The left-hand side has normal form
\[
\begin{aligned}
  N\bigl((s\otimes t)\fatsemi(u\otimes v)\bigr)
  =
  \bigl\{
    [(s'\otimes t')\fatsemi(u'\otimes v')]_0
    \mathrel{\big|}
    {}&
    [s']_0\in N(s),\
    [t']_0\in N(t),\\
    &
    [u']_0\in N(u),\
    [v']_0\in N(v)
  \bigr\}.
\end{aligned}
\]
The right-hand side has normal form
\[
\begin{aligned}
  N\bigl((s\fatsemi u)\otimes(t\fatsemi v)\bigr)
  =
  \bigl\{
    [(s'\fatsemi u')\otimes(t'\fatsemi v')]_0
    \mathrel{\big|}
    {}&
    [s']_0\in N(s),\
    [t']_0\in N(t),\\
    &
    [u']_0\in N(u),\
    [v']_0\in N(v)
  \bigr\}.
\end{aligned}
\]
These sets are equal because
\[
  [(s'\otimes t')\fatsemi(u'\otimes v')]_0
  =
  [(s'\fatsemi u')\otimes(t'\fatsemi v')]_0
\]
by bifunctoriality in \(\equiv_0\). The other symmetric monoidal
equations are treated analogously.

The semilattice equations become the corresponding equations for
union. Permutation, flattening and contraction all leave the union
\[
  N(f_1\join\cdots\join f_n)
  =
  N(f_1)\cup\cdots\cup N(f_n)
\]
unchanged: permutation reorders the summands, flattening replaces a
summand $N(g_1\join\cdots\join g_m)$ by $N(g_1)\cup\cdots\cup N(g_m)$,
and contraction removes a repeated summand.

For example, bilinearity of composition gives
\begin{align*}
  N((f_1\join\cdots\join f_n)\fatsemi h)
  &=
  \bigl\{
    [u\fatsemi v]_0
    \mathrel{\big|}
    [u]_0\in N(f_1)\cup\cdots\cup N(f_n),\
    [v]_0\in N(h)
  \bigr\}\\
  &=
  N\bigl((f_1\fatsemi h)\join\cdots\join(f_n\fatsemi h)\bigr),
\end{align*}
and the remaining bilinearity equations are analogous. Moreover,
\begin{align*}
  N\bigl(\Lambda(f_1\join\cdots\join f_n)\bigr)
  &=
  \bigl\{
    [\Lambda(u)]_0
    \mathrel{\big|}
    [u]_0\in N(f_1)\cup\cdots\cup N(f_n)
  \bigr\}\\
  &=
  N\bigl(\Lambda(f_1)\join\cdots\join\Lambda(f_n)\bigr).
\end{align*}
An induction on the derivation therefore gives
\[
  s\equiv_{\mathsf{str}}t
  \quad\Longrightarrow\quad
  N(s)=N(t).
  % \tag{\ref{lemma:lambda-term-normal-forms}.1}
  % \label{eq:lambda-term-normal-forms-1}
\]

For the converse, for every finite nonempty set \(S\) of
\(\equiv_0\)-classes, choose one representative of each class and an
enumeration, and write
\[
  \bigjoin S
\]
for the join of the chosen representatives in that order; for a
singleton this is the representative itself.  
The
\(\equiv_{\mathsf{str}}\)-class of \(\bigjoin S\) is independent of the
remaining choices: changing a representative uses \(\equiv_0\), and
changing the enumeration uses permutation.

We show by structural induction that
\begin{equation}
  t\equiv_{\mathsf{str}}\bigjoin N(t).
  \tag{$\ast$}
  \label{eq:lambda-term-normal-forms-2}
\end{equation}

If \(t\) is atomic, then \(N(t)=\{[t]_0\}\), so the result follows
immediately.

If \(t=t_1\join\cdots\join t_n\), the induction hypotheses give
\[
  t_i\equiv_{\mathsf{str}}\bigjoin N(t_i)
  \qquad(1\leq i\leq n).
\]
Hence
\[
\begin{aligned}
  t_1\join\cdots\join t_n
  &\equiv_{\mathsf{str}}
  \bigjoin N(t_1)\join\cdots\join\bigjoin N(t_n)\\
  &\equiv_{\mathsf{str}}
  \bigjoin\bigl(N(t_1)\cup\cdots\cup N(t_n)\bigr)
  =
  \bigjoin N(t),
\end{aligned}
\]
where the second step flattens the nested joins and then contracts the
classes occurring in more than one \(N(t_i)\).

If \(t=t_1\fatsemi t_2\), write
\[
  \bigjoin N(t_1)=\bigjoin_i u_i,
  \qquad
  \bigjoin N(t_2)=\bigjoin_j v_j.
\]
The induction hypotheses and bilinearity give
\[
\begin{aligned}
  t_1\fatsemi t_2
  &\equiv_{\mathsf{str}}
  \left(\bigjoin_i u_i\right)
  \fatsemi
  \left(\bigjoin_j v_j\right)\\
  &\equiv_{\mathsf{str}}
  \bigjoin_{i,j}(u_i\fatsemi v_j)\\
  &\equiv_{\mathsf{str}}
  \bigjoin N(t).
\end{aligned}
\]
Here the second step applies bilinearity in the left argument, then
bilinearity in the right argument of each resulting summand, and
finally flattens the nested joins; the third step contracts repeated
\(\equiv_0\)-classes.
The tensor case is analogous.

Finally, suppose \(t=\Lambda(u)\), and write
\[
  \bigjoin N(u)=\bigjoin_i u_i.
\]
The induction hypothesis and preservation of joins by abstraction give
\[
\begin{aligned}
  \Lambda(u)
  &\equiv_{\mathsf{str}}
  \Lambda\left(\bigjoin_i u_i\right)\\
  &\equiv_{\mathsf{str}}
  \bigjoin_i\Lambda(u_i)\\
  &\equiv_{\mathsf{str}}
  \bigjoin N(\Lambda(u)).
\end{aligned}
\]
This proves \eqref{eq:lambda-term-normal-forms-2}.

If \(N(s)=N(t)\), choose the same representative expression for their
common normal-form set. Then
\[
  s
  \equiv_{\mathsf{str}}
  \bigjoin N(s)
  =
  \bigjoin N(t)
  \equiv_{\mathsf{str}}
  t,
\]
and hence \(s\equiv_{\mathsf{str}}t\).
\end{proof}

\subsection{Complete resolutions}

For a closed extended cospan $\mathcal G$, let
\[
 \operatorname{Flat}(\mathcal G)
\]
be the set of isomorphism classes of its complete
resolutions, defined analogously to~\cref{def:flat_resolutions}. 
Each is a join-free closed extended cospan of the original
external type, still equipped with the blocks of its surviving
$\lambda$-bodies.

\begin{lemma}[Resolution and graph operations]
\label{lemma:lambda-resolution-operations}
We have the following properties of complete resolutions:
\begin{align}
 \operatorname{Flat}(\mathcal G_1\join\cdots\join\mathcal G_n)
 &=\operatorname{Flat}(\mathcal G_1)\cup\cdots\cup
   \operatorname{Flat}(\mathcal G_n),
 \label{eq:res-join}\\
 \operatorname{Flat}(\mathcal G\fatsemi\mathcal H)
 &=\{[\mathcal K\fatsemi\mathcal L]_{\cong}
       \mid[\mathcal K]_{\cong}\in\operatorname{Flat}(\mathcal G),\
            [\mathcal L]_{\cong}\in\operatorname{Flat}(\mathcal H)\},
 \label{eq:res-composition}\\
 \operatorname{Flat}(\mathcal G\otimes\mathcal H)
 &=\{[\mathcal K\otimes\mathcal L]_{\cong}
       \mid[\mathcal K]_{\cong}\in\operatorname{Flat}(\mathcal G),\
            [\mathcal L]_{\cong}\in\operatorname{Flat}(\mathcal H)\},
 \label{eq:res-tensor}\\
 \operatorname{Flat}(\Lambda(\mathcal G))
 &=\{[\Lambda(\mathcal K)]_{\cong}
       \mid[\mathcal K]_{\cong}\in\operatorname{Flat}(\mathcal G)\}.
 \label{eq:res-lambda}
\end{align}
\end{lemma}
\begin{proof}
Neither composition nor tensor introduce new $\boxsort$-edges, so the sets of complete resolutions are given by those of the components, combined using composition or tensor respectively.
Under $\join$ the $\boxsort$-edges of $\mathcal{G}_1, \ldots, \mathcal{G}_n$ end up in different components of the fresh $\boxsort$-edge, and a complete resolution selects one of them, so the set of complete resolutions is the union of the $n$ sets.
Finally, $\lambdasort$-edges are not removed, so the complete resolutions of $\Lambda(\mathcal{G})$ are given by those of $\mathcal{G}$ with the $\lambdasort$-edge and its root retained.
\end{proof}

\begin{lemma}[Resolution of an interpreted term]
\label{lemma:lambda-resolution-interpretation}
For every raw $\Sigma^{\join}$-term $t \in \operatorname{Tm}^{\Lambda,\join}_{\Sigma}(A,B)$,
\[
 \operatorname{Flat}(\widehat t)
 =\{[\widehat u]_{\cong}\mid[u]_0\in N(t)\}.
\]
\end{lemma}
\begin{proof}
The set on the right is well-defined: if $[u]_0 = [v]_0$, then $\widehat{u} \cong \widehat{v}$.
We proceed by structural induction on $t$.
Each atomic term has its own interpretation
as its sole resolution.
Consider the case $t = t_1 \fatsemi t_2$.
By the induction hypothesis, we have
\[
\operatorname{Flat}(\widehat{t_1}) = \{[\widehat{u_1}]_{\cong} \mid [u_1]_0 \in N(t_1)\}
\]
and
\[
\operatorname{Flat}(\widehat{t_2}) = \{[\widehat{u_2}]_{\cong} \mid [u_2]_0 \in N(t_2)\}~.
\]
By~\cref{lemma:lambda-resolution-operations}, we have 
\begin{align*}
  \operatorname{Flat}(\widehat t)
  &=
  \operatorname{Flat}
  (\widehat{t_1}\fatsemi\widehat{t_2})\\
  &=
  \bigl\{
    [\mathcal K\fatsemi\mathcal L]_{\cong}
    \mathrel{\big|}
    [\mathcal K]_{\cong}
      \in\operatorname{Flat}(\widehat{t_1}),
    \
    [\mathcal L]_{\cong}
      \in\operatorname{Flat}(\widehat{t_2})
  \bigr\}\\
  &=
  \bigl\{
    [\widehat{u_1}\fatsemi\widehat{u_2}]_{\cong}
    \mathrel{\big|}
    [u_1]_0\in N(t_1),
    \
    [u_2]_0\in N(t_2)
  \bigr\}\\
  &=
  \bigl\{
    [\widehat{u_1\fatsemi u_2}]_{\cong}
    \mathrel{\big|}
    [u_1]_0\in N(t_1),
    \
    [u_2]_0\in N(t_2)
  \bigr\}\\
  &=
  \{[\widehat u]_{\cong}\mid[u]_0\in N(t)\}.
\end{align*}
For the join case $t=t_1\join\cdots\join t_n$, the induction hypothesis
and \eqref{eq:res-join} give
\[
  \operatorname{Flat}(\widehat t)
  =
  \bigcup_{i=1}^{n}\operatorname{Flat}(\widehat{t_i})
  =
  \bigcup_{i=1}^{n}
  \{[\widehat u]_{\cong}\mid[u]_0\in N(t_i)\}
  =
  \{[\widehat u]_{\cong}\mid[u]_0\in N(t)\},
\]
by the definition of $N$ on an $n$-ary join. The tensor and abstraction
cases are analogous to the case of composition, using
\eqref{eq:res-tensor} and \eqref{eq:res-lambda}.
\end{proof}

\begin{lemma}[Rewriting congruence]
\label{lemma:lambda-rewriting-congruence}
Let $\mathcal{G}, \mathcal{L}, \mathcal{R}, \mathcal{H}$ be well-typed
MDA closed extended cospans and suppose that
\[
\mathcal{G} \widetilde{\rightsquigarrow}_{\langle \mathcal{L}, \mathcal{R} \rangle} \mathcal{H}
\]
is a convex closed EDPOI rewrite at a convex match
$m : \mathcal{L} \to \mathcal{G}$.
Then, whenever the displayed composites are defined, we have (1)-(3) as
in~\cref{lemma:rewriting_congruence} and furthermore,
\[
\Lambda(\mathcal{G}) \widetilde{\rightsquigarrow}_{\langle \mathcal{L}, \mathcal{R} \rangle} \Lambda(\mathcal{H}).
\]
In each case the match of the resulting rewrite is the image of $m$
under the corresponding operation, and its extended boundary complement
extends that of $m$ by the fixed operands.
\end{lemma}
\begin{proof}
Analogous to~\cref{lemma:rewriting_congruence}.
\end{proof}

\begin{lemma}[Resolution locality]
\label{lemma:lambda-resolution-locality}
Suppose $\mathcal L,\mathcal R:\partial A\to\partial B$ have the same
external type and
$\operatorname{Flat}(\mathcal L)=\operatorname{Flat}(\mathcal R)$.
If a convex closed EDPOI rewrite
$\mathcal G\mathrel{\widetilde{\rightsquigarrow}_{\langle\mathcal L,\mathcal R\rangle}}\mathcal H$
replaces $\mathcal L$ by $\mathcal R$, then
\[
 \operatorname{Flat}(\mathcal G)=\operatorname{Flat}(\mathcal H).
\]
The strictly internal interfaces of $\mathcal G$ and $\mathcal H$ need
not coincide before resolution.
\end{lemma}
\begin{proof}
Since the match $m : \mathcal{L} \to \mathcal{G}$ is convex, there is a
decomposition of $\mathcal{G}$ as in~\cref{lemma:lambda-convex-decomposition}.
Under such a decomposition, $\mathcal{G}$ is obtained by repeated
composition, tensor, join and currying around $\mathcal{L}$, and the
decomposition isomorphism carries the canonical occurrence of
$\mathcal L$ to $m$. Starting from the identity application
$\mathcal L\mathrel{\widetilde{\rightsquigarrow}_{\langle\mathcal L,\mathcal R\rangle}}\mathcal R$
(as in~\cref{lemma:edpoi_identity_application}) and applying
\cref{lemma:lambda-rewriting-congruence} to each of these operations
yields a rewrite at the canonical occurrence whose result is the same
construction around $\mathcal R$. Its extended boundary complement is
transported to one for $m$, so by
\cref{prop:lambda-boundary-complement-unique} the result is isomorphic
to $\mathcal H$. Hence $\mathcal H$ is obtained from $\mathcal R$ by the
same composition, tensor, join and currying operations.
Since $\operatorname{Flat}(\mathcal L) = \operatorname{Flat}(\mathcal R)$,
\cref{lemma:lambda-resolution-operations}, applied to each of these
operations, gives $\operatorname{Flat}(\mathcal{G}) = \operatorname{Flat}(\mathcal{H})$.
\end{proof}

Let $\mathfrak{R}_{\Lambda}$ be the set of structural rules obtained as in~\cref{def:structural_rules}, additionally with the following schema corresponding to the equation~\eqref{eq:lambda-join-structural}:
\[
\adjustbox{max width=\textwidth}{
\tikzfig{./figures/lambda/semilat_distributivity_law1}~.
}
\]

Write $\mathcal{G} \sim_{\mathfrak{R}_{\Lambda}} \mathcal{H}$ if $\mathcal{G} \widetilde{\leftrightsquigarrow}_{\mathfrak{R}_{\Lambda}}^{*} \mathcal{H}$, the reflexive, symmetric and transitive closure of $\widetilde{\rightsquigarrow}_{\mathfrak{R}_{\Lambda}}$ of~\cref{def:lambda-convex-edpoi-rewriting}.
By~\cref{lemma:lambda-rewriting-congruence}, $\sim_{\mathfrak{R}_{\Lambda}}$ is a congruence.

\begin{corollary}[Invariance of resolution]
\label{cor:lambda-resolution-invariance}
If $\mathcal G\sim_{\mathfrak R_\Lambda}\mathcal H$, then
$\operatorname{Flat}(\mathcal G)=\operatorname{Flat}(\mathcal H)$.
\end{corollary}
\begin{proof}
This follows from~\cref{lemma:lambda-resolution-locality} since every structural rewrite rule preserves $\operatorname{Flat}$.
\end{proof}

For a finite nonempty set $S$ of isomorphism classes of join-free closed
extended cospans of one external type, choose a representative of each
class and put
\[
 \operatorname{Can}(S)=\bigjoin_{[\mathcal K]\in S}\mathcal K,
\]
the $n$-ary join of the chosen representatives. For a singleton this is
the chosen representative itself. Since the join is $n$-ary, no
bracketing has to be chosen, and the components of a $\boxsort$-edge
carry no order, so $\operatorname{Can}(S)$ is independent of the
enumeration up to isomorphism
(\cref{remark:commutativity_absorbed}); different choices of
representatives give isomorphic results.


\begin{lemma}[Graphical normal form]
\label{lemma:lambda-graphical-normal-form}
For every raw term $t$,
\[
  \widehat t
  \sim_{\mathfrak R_{\Lambda}}
  \operatorname{Can}
  \bigl(\operatorname{Flat}(\widehat t)\bigr).
\]
\end{lemma}
\begin{proof}
% We proceed by structural induction on \(t\). 
% If \(t\) is atomic, then
% \[
%   \operatorname{Flat}(\widehat t)
%   =
%   \{[\widehat t]_{\cong}\},
% \]
% so the claim follows immediately from the definition of
% \(\operatorname{Can}\).

% Suppose \(t=t_1\join\cdots\join t_n\). 
% By the induction hypotheses and
% closure of \(\sim_{\mathfrak R_\Lambda}\) under join,
% \[
%   \widehat{t_1}\join\cdots\join\widehat{t_n}
%   \sim_{\mathfrak R_\Lambda}
%   \operatorname{Can}(\operatorname{Flat}(\widehat{t_1}))
%   \join\cdots\join
%   \operatorname{Can}(\operatorname{Flat}(\widehat{t_n})).
% \]
% Flattening the nested joins and contracting repeated isomorphism
% classes turns the right-hand side into
% \[
%   \operatorname{Can}\bigl(
%     \operatorname{Flat}(\widehat{t_1})
%     \cup\cdots\cup
%     \operatorname{Flat}(\widehat{t_n})
%   \bigr),
% \]
% which is
% \[
%   \operatorname{Can}(\operatorname{Flat}(\widehat t))
% \]
% by \cref{eq:res-join}.

% Suppose \(t=t_1\fatsemi t_2\). 
% Using the induction hypotheses and closure under composition, we obtain
% \[
%   \widehat{t_1}\fatsemi\widehat{t_2}
%   \sim_{\mathfrak R_\Lambda}
%   \operatorname{Can}(\operatorname{Flat}(\widehat{t_1}))
%   \fatsemi
%   \operatorname{Can}(\operatorname{Flat}(\widehat{t_2})).
% \]
% Bilinearity in the left argument, then in the right argument of each
% resulting summand, followed by flattening, expresses the right-hand
% side as the join of all cospans
% \[
%   \mathcal K\fatsemi\mathcal L,
%   \qquad
%   [\mathcal K]_{\cong}
%     \in\operatorname{Flat}(\widehat{t_1}),
%   \qquad
%   [\mathcal L]_{\cong}
%     \in\operatorname{Flat}(\widehat{t_2}).
% \]
% Contraction removes repeated isomorphism classes. By
% \cref{eq:res-composition}, the result is structurally equivalent to
% \[
%   \operatorname{Can}(\operatorname{Flat}(\widehat t)).
% \]
% The tensor case is analogous, using \cref{eq:res-tensor}.

% Finally, suppose \(t=\Lambda(u)\). 
% By the induction hypothesis and because $\sim_{\mathfrak{R}_\Lambda}$ is a congruence,
% \[
%   \Lambda(\widehat u)
%   \sim_{\mathfrak R_\Lambda}
%   \Lambda\bigl(
%     \operatorname{Can}(\operatorname{Flat}(\widehat u))
%   \bigr).
% \]
% The structural rule
% \[
%   \Lambda(\mathcal G_1\join\cdots\join\mathcal G_n)
%   \sim_{\mathfrak R_\Lambda}
%   \Lambda(\mathcal G_1)\join\cdots\join\Lambda(\mathcal G_n)
% \]
% moves the join outside the abstraction. 
% Flattening, permutation and
% contraction then give
% \[
%   \Lambda\bigl(
%     \operatorname{Can}(\operatorname{Flat}(\widehat u))
%   \bigr)
%   \sim_{\mathfrak R_\Lambda}
%   \operatorname{Can}
%   \left(
%     \{
%       [\Lambda(\mathcal K)]_{\cong}
%       \mid
%       [\mathcal K]_{\cong}
%         \in\operatorname{Flat}(\widehat u)
%     \}
%   \right).
% \]
% By \cref{eq:res-lambda}, the latter is
% \[
%   \operatorname{Can}(\operatorname{Flat}(\widehat t)).
% \]
As in~\cref{lemma:canonical_representative}, we can turn $\widehat{t}$
into the form $\bigjoin_i \mathcal{F}_{i}$ such that each $\mathcal{F}_{i}$
is box-free and $\mathcal{F}_{i} \not \cong \mathcal{F}_{j}$ for
$i \neq j$. The only additional step compared with that lemma is that
the schema for~\eqref{eq:lambda-join-structural} moves joins out of
$\lambdasort$-bodies, so that no $\boxsort$-edge remains below a
$\lambdasort$-edge. If $\widehat t$ is already box-free, the join has a
single summand, which by our convention is $\mathcal F_1$ itself.
Since $\operatorname{Flat}$ is invariant under $\sim_{\mathfrak{R}_{\Lambda}}$
(\cref{cor:lambda-resolution-invariance}) and since each
$\mathcal{F}_{i}$ is box-free, we have
\[
\operatorname{Flat}\Bigl(\bigjoin_i \mathcal{F}_{i}\Bigr) = \bigcup_{i} \operatorname{Flat}(\mathcal{F}_{i}) = \bigcup_{i} \{[\mathcal{F}_{i}]_{\cong}\} = \operatorname{Flat}(\widehat{t}),
\]
and therefore
\[
\widehat{t} \sim_{\mathfrak{R}_{\Lambda}} \bigjoin_i \mathcal{F}_{i} \cong \operatorname{Can}(\operatorname{Flat}(\widehat{t})).
\]
\end{proof}


\subsection{The adequacy theorem}

\begin{theorem}[Structural adequacy]
\label{theorem:lambda-structural-adequacy}
For raw closed $\Sigma^{\join}$-terms $s,t \in \operatorname{Tm}^{\Lambda, \join}_{\Sigma}(A, B)$,
\[
 s\equiv_{\mathsf{str}}t
 \quad\Longleftrightarrow\quad
 \widehat s\sim_{\mathfrak R_{\Lambda}}\widehat t.
\]
\end{theorem}
\begin{proof}
Suppose first that $s\equiv_{\mathsf{str}}t$.  By
\cref{lemma:lambda-term-normal-forms}, $N(s)=N(t)$, and hence
\cref{lemma:lambda-resolution-interpretation} gives
\[
 \operatorname{Flat}(\widehat s)
 =\operatorname{Flat}(\widehat t).
\]
Applying~\cref{lemma:lambda-graphical-normal-form} to both terms gives
\[
 \widehat s
 \sim_{\mathfrak R_{\Lambda}}
 \operatorname{Can}(\operatorname{Flat}(\widehat s))
 \sim_{\mathfrak R_{\Lambda}}
 \operatorname{Can}(\operatorname{Flat}(\widehat t))
 \sim_{\mathfrak R_{\Lambda}}
 \widehat t.
\]

Conversely, suppose
$\widehat s\sim_{\mathfrak R_{\Lambda}}\widehat t$. 
By~\cref{cor:lambda-resolution-invariance}, structural rules preserve
the set of complete resolutions, so
\[
 \operatorname{Flat}(\widehat s)
 =\operatorname{Flat}(\widehat t).
\]
Using~\cref{lemma:lambda-resolution-interpretation}, this says
\[
 \bigl\{[\widehat u]_{\cong}\mid[u]_{0}\in N(s)\bigr\}
 =
 \bigl\{[\widehat v]_{\cong}\mid[v]_{0}\in N(t)\bigr\}.
\]
The map $[u]_{0}\mapsto[\widehat u]_{\cong}$ is injective by
\cref{prop:lambda-join-free-faithfulness}.  Therefore $N(s)=N(t)$, and
\cref{lemma:lambda-term-normal-forms} yields
$s\equiv_{\mathsf{str}}t$.
\end{proof}

\section{The rewriting correspondence}
\label{sec:lambda-rewriting-correspondence}

We now compare directed rewriting of raw closed-monoidal terms modulo
$\equiv_{\mathsf{str}}$ with directed EDPOI rewriting modulo
$\mathfrak R_{\Lambda}$.  Throughout this section a rewrite rule is a pair
\[
 \rho=\langle l, r\rangle,
 \qquad
 l,r\in\operatorname{Tm}^{\Lambda,\join}_{\Sigma}(P,Q),
\]
of parallel raw closed-monoidal terms. 
\subsection{Rewriting modulo structural equality}

\begin{definition}[Typed one-hole contexts]
\label{def:lambda-one-hole-context}
Typed one-hole contexts with a hole of type $P\to Q$ are generated by the
grammar
\begin{equation}
\begin{aligned}
  C[-] \Coloneqq{}
  &[-]
  \mid C[-]\fatsemi t
  \mid t\fatsemi C[-]
  \mid C[-]\otimes t
  \mid t\otimes C[-]\\
  &\mid t_1\join\cdots\join C[-]\join\cdots\join t_n
  \mid \Lambda_{A,X,B}(C[-]),
\end{aligned}
\label{eq:typed-one-hole-context}
\end{equation}
where $t$ and the $t_i$ range over raw closed-monoidal terms of
$\operatorname{Tm}^{\Lambda,\join}_{\Sigma}$, $n\geq2$, $X$ is atomic, and
every clause is required to be well-typed. For a raw term
$s\in\operatorname{Tm}^{\Lambda,\join}_{\Sigma}(P,Q)$ we write $C[s]$ for the
raw term obtained by substituting $s$ for the hole; if $C[-]$ has outer
type $X\to Y$, then $C[s]\in\operatorname{Tm}^{\Lambda,\join}_{\Sigma}(X,Y)$.
\end{definition}

\begin{definition}[Term rewriting modulo structural equality]
\label{def:lambda-term-rewriting-modulo}
For a rewrite rule $\langle l, r \rangle$ with
$l,r \in \operatorname{Tm}^{\Lambda, \join}_{\Sigma}(P,Q)$ and raw closed-monoidal terms
$a,b \in \operatorname{Tm}^{\Lambda, \join}_{\Sigma}(X,Y)$, write
\[
 a\longrightarrow_{\langle l, r \rangle/\mathsf{str}}b
\]
if there are raw closed-monoidal terms $a',b'$ and a typed one-hole context $C[-]$ with a
hole of type $P\to Q$ and outer type $X\to Y$ such that
\[
 a\equiv_{\mathsf{str}}a'=C[l],
 \qquad
 b'=C[r]\equiv_{\mathsf{str}}b.
 \tag{\(\dagger\)}
 \label{eq:lambda-term-step-modulo}
\]
\end{definition}


\begin{definition}[EDPOI rewriting modulo structural rules]
\label{def:lambda-edpoi-rewriting-modulo}
For closed extended cospans $\mathcal G$ and $\mathcal H$, and a rewrite rule $\langle \mathcal{L}, \mathcal{R} \rangle$, write
\[
 \mathcal G
 \mathrel{\widetilde{\rightsquigarrow}^{\,\mathfrak R_{\Lambda}}_{\langle \mathcal{L}, \mathcal{R} \rangle}}
 \mathcal H
\]
if there are $\mathcal G'$ and $\mathcal H'$ such that
\[
 \mathcal G\sim_{\mathfrak R_{\Lambda}}\mathcal G',
 \qquad
 \mathcal G'
 \mathrel{\widetilde{\rightsquigarrow}_{\langle \mathcal{L}, \mathcal{R} \rangle}}
 \mathcal H',
 \qquad
 \mathcal H'\sim_{\mathfrak R_{\Lambda}}\mathcal H,
 \tag{\(\ddagger\)}
 \label{eq:lambda-graph-step-modulo}
\]
where the middle relation is a convex closed EDPOI application in the sense of
\cref{def:lambda-convex-edpoi-rewriting}.
\end{definition}

\subsection{Reification, context extraction and replacement}

\begin{lemma}[Reification]
\label{lemma:lambda-reification}
Let
\[
 \mathbf G
 =
 \left(
 n
 \xrightarrow{f_{\mathrm{ext}}}
 n'
 \xrightarrow{f_{\mathrm{int}}}
 \mathcal G
 \xleftarrow{g_{\mathrm{int}}}
 k'
 \xleftarrow{g_{\mathrm{ext}}}
 k
 \right)
\]
be a well-typed MDA closed extended cospan, and put
\[
 A\defeq\left\lVert\operatorname{ty}_{n}([n])\right\rVert_{\otimes},
 \qquad
 B\defeq\left\lVert\operatorname{ty}_{k}([k])\right\rVert_{\otimes},
\]
so that $n=\partial A$ and $k=\partial B$.  Then there is a raw term
$u\in\operatorname{Tm}^{\Lambda,\join}_{\Sigma}(A,B)$ such that
\[
 \widehat u\cong\mathbf G.
\]
\end{lemma}
\begin{proof}
We proceed by induction on the number $\nu(\mathbf G)$ of edges of
$\mathcal G$ which are not component edges, that is, of ordinary,
$\boxsort$- and $\lambdasort$-edges.  

Suppose first that $\nu(\mathbf G)=0$, then the only edge
of $\mathcal G$ is $\top_{\mathcal G}$, and every vertex has immediate
parent $\top_{\mathcal G}$. Moreover, the strictly internal interfaces are empty.
Every vertex has in- and out-degree zero, so the
MDA conditions make both external legs bijections onto
$V_{\mathcal G}$.  Hence $\mathbf G$ is the interpretation of a term
built from primitive identities and symmetries by $\fatsemi$ and
$\otimes$, exactly as in the base case of the proof of
\cref{theorem:monogamous_isomorphism}.

Suppose now that $\nu(\mathbf G)>0$.  
Then we may choose a non-component edge $e$ with $\top_{\mathcal{G}} <^{\mu}_{\mathcal{G}} e$.
Put
\[
 P\defeq\left\lVert l_V^*(s(e))\right\rVert_{\otimes},
 \qquad
 Q\defeq\left\lVert l_V^*(t(e))\right\rVert_{\otimes}.
\]
Let $\mathcal L$ be the sub-e-hypergraph of $\mathcal G$ consisting of
$e$, the vertices of $s(e)$ and $t(e)$, and all descendants of $e$,
with a fresh root whose immediate children are $e$ and these ports.
Its ordered input and output ports form the external interfaces
$\partial P$ and $\partial Q$.  The interface blocks of
$\mathbf G$ belonging to components below $e$ are its strictly
internal interfaces.  We obtain a well-typed MDA closed extended
cospan
\[
 \mathbf L:\partial P\longrightarrow\partial Q.
\]
Now consider the following cases.

\begin{enumerate}
 \item If $e$ is an ordinary edge labelled by $g$, then $e$ has no
 descendants and $\mathbf L\cong\widehat g$.  Put $w\defeq g$.

 \item If $e$ is a $\boxsort$-edge, let $c_1,\ldots,c_s$ be its
 components.
Each $c_i$ together with its
 descendants and its interface blocks determines a well-typed MDA
 closed extended cospan $\mathbf G_{c_i}$. 
 Well-typedness of $e$ gives
 \[
  \mathbf G_{c_i}:\partial P\longrightarrow\partial Q,
 \]
 and
 \[
  \mathbf L\cong
  \mathbf G_{c_1}\join\cdots\join\mathbf G_{c_s}.
 \]
 Since $\nu(\mathbf G_{c_i})<\nu(\mathbf G)$, the induction hypothesis
 gives raw closed-monoidal terms $u_i:P\to Q$ with
 $\widehat{u_i}\cong\mathbf G_{c_i}$.  
 Put
 \[
  w\defeq u_1\join\cdots\join u_s.
 \]
 This is the step where we use $n$-ary joins in the syntax; otherwise we would have to reify terms only up to structural equivalence.

 \item If $e$ is a $\lambdasort$-edge with unique component $c$,
 well-typedness of $e$ gives objects $X,C$ and an atomic type $D$ with
 \[
  P=D,
  \qquad
  Q=X\multimap C,
  \qquad
  \mathbf G_c:\partial(D \otimes X)\longrightarrow\partial C.
 \]
 Therefore,
 \[
  \mathbf L\cong\Lambda(\mathbf G_c).
 \]
%  This is the identification used in case~(2b) of
%  \cref{lemma:lambda-convex-decomposition}.  
 Since
 $\nu(\mathbf G_c)<\nu(\mathbf G)$, the induction hypothesis gives a
 raw term $v:D\otimes X\to C$ with $\widehat v\cong\mathbf G_c$.  Put
 \[
  w\defeq\Lambda_{D,X,C}(v).
 \]
\end{enumerate}

Hence, we get a raw term $w : P \to Q$ such that
\[
 \widehat w\cong\mathbf L.
\]

The inclusion $m:\mathcal L\to\mathcal G$ is an anchored monomorphism
with anchor $\top_{\mathcal G}$.  Its image contains every descendant
of $e$, and the non-root edges of $\mathcal L$ are $e$ and its
descendants.  Hence the image is closed under descendants below every
matched non-root edge.  The image is also convex. 
Case~(1) of
\cref{lemma:lambda-convex-decomposition} therefore gives a typed
interface $q$ and well-typed MDA closed extended cospans
\[
 \mathbf C_1:\partial A\longrightarrow q\coprod\partial P,
 \qquad
 \mathbf C_2:q\coprod\partial Q\longrightarrow\partial B
\]
such that
\[
 \mathbf G
 \cong
 \mathbf C_1\fatsemi(\mathbf 1_q\otimes\mathbf L)\fatsemi\mathbf C_2.
\]
Composition and tensor glue only along interface vertices, so the
non-component edges of $\mathcal G$ are partitioned between
$\mathbf C_1$, $\mathbf L$ and $\mathbf C_2$.  Since $e$ lies in
$\mathbf L$,
\[
 \nu(\mathbf C_1)<\nu(\mathbf G)
 \qquad\text{and}\qquad
 \nu(\mathbf C_2)<\nu(\mathbf G).
\]
The induction hypothesis gives raw closed-monidal terms $c_1$ and $c_2$ with
$\widehat{c_1}\cong\mathbf C_1$ and $\widehat{c_2}\cong\mathbf C_2$.
Write $K\defeq\lVert\operatorname{ty}_q([\,|q|\,])\rVert_{\otimes}$, so that $\partial K=q$, and put
\[
 u\defeq c_1\fatsemi(\id_K\otimes w)\fatsemi c_2 .
\]
Composition and tensor preserve isomorphisms of closed extended
cospans.  By compositionality of the interpretation,
\[
 \widehat u
 =
 \widehat{c_1}\fatsemi(\mathbf 1_{q}\otimes\widehat w)\fatsemi\widehat{c_2}
 \cong
 \mathbf C_1\fatsemi(\mathbf 1_q\otimes\mathbf L)\fatsemi\mathbf C_2
 \cong
 \mathbf G.
\]
\end{proof}

This lemma allows us to turn a well-typed MDA closed extended cospan into a raw closed-monoidal term, which we use to reconstruct the term context from a convex match in what follows.
% In other words, the interpretation of
% \cref{def:lambda-concrete-interpretation} is surjective on isomorphism
% classes of well-typed MDA closed extended cospans.  Combined with
% \cref{theorem:lambda-structural-adequacy}, it yields, for all objects
% $A,B$, a bijection
% \[
%  \operatorname{Tm}^{\Lambda,\join}_{\Sigma}(A,B)/{\equiv_{\mathsf{str}}}
%  \;\cong\;
%  \catname{CMEHypI}_{\bullet}(\Sigma)(\partial A,\partial B)/{\sim_{\mathfrak R_{\Lambda}}}.
% \]
% By compositionality of the interpretation these bijections respect
% composition, tensor, join and abstraction, and $\partial$ is a bijection
% on objects, since $\operatorname{wd}$ and $\lVert-\rVert_\otimes$ are
% mutually inverse. They therefore assemble into an isomorphism of
% $\catname{SLat}$-enriched symmetric monoidal categories
% \[
%  \operatorname{Tm}^{\Lambda,\join}_{\Sigma}/{\equiv_{\mathsf{str}}}
%  \;\cong\;
%  \catname{CMEHypI}_{\bullet}(\Sigma)/\mathfrak R_{\Lambda}.
% \]
% This is the $\Lambda$-analogue of
% \cref{theorem:mehypi_slat_isomorphism}.

\subsection{Soundness and completeness}

\begin{theorem}[Rewriting correspondence]
\label{theorem:lambda-rewriting-correspondence}
Let $l, r \in \operatorname{Tm}^{\Lambda,\join}_{\Sigma}(P,Q)$ be raw
terms, and let $\rho = \langle l, r \rangle$ be the corresponding
rewrite rule. For raw closed-monoidal terms
$a, b \in \operatorname{Tm}^{\Lambda,\join}_{\Sigma}(X,Y)$,
\[
 a \longrightarrow_{\rho/\mathsf{str}} b
 \quad\Longleftrightarrow\quad
 \widehat a
 \mathrel{\widetilde{\rightsquigarrow}^{\,\mathfrak R_{\Lambda}}_{\langle\widehat l,\widehat r\rangle}}
 \widehat b.
\]
\end{theorem}
\begin{proof}
This follows the same pattern as the proof
of~\cref{theorem:convex_edpoi_rewriting_slat}, with the difference that
reification (\cref{lemma:lambda-reification}) takes the place of the
quotient isomorphism used there.

For the forward direction, we first record the following fact.
For every typed one-hole context \(C[-]\), the canonical occurrence
\[
  \iota_l:
  \widehat l\longrightarrow\widehat{C[l]}
\]
is a convex match, and replacing \(\widehat l\) by \(\widehat r\) using
the corresponding complement gives
\[
  \widehat{C[l]}
  \mathrel{\widetilde{\rightsquigarrow}_{\langle\widehat l,\widehat r\rangle}}
  \widehat{C[r]}.
  \tag{\(\ast\)}
  \label{eq:lambda-context-replacement-in-proof}
\]
This is proved by induction on \(C[-]\).

For \(C[-]=[-]\), we have \(C[l]=l\) and \(C[r]=r\), so as
in~\cref{lemma:edpoi_identity_application} we immediately get
\[
  \widehat l
  \mathrel{\widetilde{\rightsquigarrow}_{\langle\widehat l,\widehat r\rangle}}
  \widehat r .
\]
Furthermore, the identity match is convex.

For the inductive step, a typed one-hole context is of one of the
forms~\eqref{eq:typed-one-hole-context} of \cref{def:lambda-one-hole-context}. Say
\(C[-]=C'[-]\fatsemi t\). The induction hypothesis gives a convex closed
EDPOI rewrite
\[
  \widehat{C'[l]}
  \mathrel{\widetilde{\rightsquigarrow}_{\langle\widehat l,\widehat r\rangle}}
  \widehat{C'[r]}.
\]
Clause~(1) of~\cref{lemma:lambda-rewriting-congruence} gives
\[
  \widehat{C'[l]}\fatsemi\widehat t
  \mathrel{\widetilde{\rightsquigarrow}_{\langle\widehat l,\widehat r\rangle}}
  \widehat{C'[r]}\fatsemi\widehat t,
\]
and by compositionality of the interpretation
(\cref{def:lambda-concrete-interpretation}) its source and target are
\(\widehat{C[l]}\) and \(\widehat{C[r]}\). The other cases follow by
applying the corresponding clauses
of~\cref{lemma:lambda-rewriting-congruence}.

Now suppose that
\[
  a\longrightarrow_{\rho/\mathsf{str}}b.
\]
By \cref{def:lambda-term-rewriting-modulo}, there is a typed one-hole
context \(C[-]\) such that
\[
  a\equiv_{\mathsf{str}}C[l],
  \qquad
  C[r]\equiv_{\mathsf{str}}b.
\]
By \eqref{eq:lambda-context-replacement-in-proof}, there is a direct
convex closed EDPOI rewrite
\[
  \widehat{C[l]}
  \mathrel{\widetilde{\rightsquigarrow}_{\langle\widehat l,\widehat r\rangle}}
  \widehat{C[r]}.
\]
Structural adequacy (\cref{theorem:lambda-structural-adequacy}) gives
\[
  \widehat a
  \sim_{\mathfrak R_\Lambda}
  \widehat{C[l]}
  \qquad\text{and}\qquad
  \widehat{C[r]}
  \sim_{\mathfrak R_\Lambda}
  \widehat b.
\]
Hence
\[
  \widehat a
  \mathrel{\widetilde{\rightsquigarrow}^{\,\mathfrak R_\Lambda}_{\langle\widehat l,\widehat r\rangle}}
  \widehat b
\]
by \cref{def:lambda-edpoi-rewriting-modulo}.

Conversely, suppose
\[
  \widehat a
  \mathrel{\widetilde{\rightsquigarrow}^{\,\mathfrak R_\Lambda}_{\langle\widehat l,\widehat r\rangle}}
  \widehat b.
\]
Choose closed extended cospans \(\mathcal G'\) and \(\mathcal H'\) such
that
\[
  \widehat a\sim_{\mathfrak R_\Lambda}\mathcal G',
  \qquad
  \mathcal G'
  \mathrel{\widetilde{\rightsquigarrow}_{\langle\widehat l,\widehat r\rangle}}
  \mathcal H',
  \qquad
  \mathcal H'\sim_{\mathfrak R_\Lambda}\widehat b.
  % \tag{\(\dagger\)}
  % \label{eq:lambda-one-step-witnesses}
\]
Let
\[
  m:\widehat l\longrightarrow\mathcal G'
\]
be the match of the middle rewrite. Since \(m\) is a convex match,
\cref{lemma:lambda-convex-decomposition} supplies a well-typed MDA
closed extended cospan \(\mathbf D\), constructed from \(\widehat l\)
and fixed well-typed MDA closed extended cospans by composition, tensor
product, join and abstraction.

For every fixed MDA closed extended cospan occurring in this finite
construction of \(\mathbf D\), choose, by
\cref{lemma:lambda-reification}, a raw term whose interpretation is
isomorphic to it. Replacing the distinguished occurrence of
\(\widehat l\) by a hole and all the fixed operands by their chosen
terms gives a typed one-hole context \(C[-]\). The identity match gives
us
\[
  \widehat l
  \mathrel{\widetilde{\rightsquigarrow}_{\langle\widehat l,\widehat r\rangle}}
  \widehat r,
\]
and by recursively lifting this rewrite through the decomposition,
as in~\eqref{eq:lambda-context-replacement-in-proof}, we get
\[
  \widehat{C[l]}
  \mathrel{\widetilde{\rightsquigarrow}_{\langle\widehat l,\widehat r\rangle}}
  \widehat{C[r]}.
\]
Since the match \(m\) factors through the isomorphism induced by the
decomposition, and because the extended boundary complement is unique
(\cref{prop:lambda-boundary-complement-unique}),
we get \(\widehat{C[r]}\cong\mathcal H'\). The decomposition itself
gives us \(\widehat{C[l]}\cong\mathcal G'\). Therefore, since we have
\[
  \widehat a\sim_{\mathfrak R_\Lambda}\mathcal G'\cong\widehat{C[l]}
  \qquad\text{and}\qquad
  \widehat b\sim_{\mathfrak R_\Lambda}\mathcal H'\cong\widehat{C[r]},
\]
by structural adequacy we get that
\[
  a\equiv_{\mathsf{str}}C[l],
  \qquad
  C[r]\equiv_{\mathsf{str}}b,
\]
and hence \(a\longrightarrow_{\rho/\mathsf{str}}b\) by
\cref{def:lambda-term-rewriting-modulo}.
\end{proof}

We thus established that a term rewrite modulo structural equality
exists exactly when the interpreted closed e-hypergraph admits the
corresponding convex EDPOI rewrite modulo structural rules. The
graphical representation absorbs the symmetric monoidal equations and
the permutation of join alternatives, while the remaining semilattice
equations are handled by $\mathfrak R_{\Lambda}$. As in
\cref{sec:monoidal_applications}, algorithmically, matching may be performed against the
hierarchical representation by following its interfaces, without first
unfolding every hierarchical edge.

\section{Categorical foundations of slotted e-graphs}
We finish with a brief discussion of the equation~\eqref{eq:lambda-join-structural} with respect to slotted e-graphs of Schneider et al.~\cite{slotted-egraphs}.
Slotted e-graphs extend the traditional signature employed by e-graphs with an explicit construct for binding ($\text{bind}\; \$x\; tc$) and variables (\emph{slots}).
	Their syntax is generated by the following grammar:
	\begin{alignat*}{2}
		\text{(functional symbols)} &        &                            & \; f,\;g, \ldots ;                                                                       \\
		\text{(slots)}              &        &                            & \; \$x,\; \$y, \ldots ;                                                                   \\
		\text{(term children)}      &        & tc               \Coloneqq & \; t \mid tc_{1},\; \ldots,\; tc_{n} \mid \text{bind}\; \$x\; tc \mid \$x ; \qquad n \geq 1 \\
		\text{(terms)}              & \qquad & t                \Coloneqq & \; f \mid f(tc) .
	\end{alignat*}
	The traditional congruence relation maintained by e-graphs for a set of equations $E$ is given below (plus the obvious rules for transitivity, symmetry and reflexivity of $\cong$; in this subsection only, $\cong$ denotes the e-graph congruence rather than isomorphism):
	% \begin{mathpar}
	% 	\inferrule*[right=start]
	% 	{a = b \in E}
	% 	{E \vdash a \cong b}
	% 	\quad
	% 	\inferrule*[right=cong. (variadic)]
	% 	{E \vdash tc_i \cong tc'_i,\; i = 1, \ldots, k}
	% 	{E \vdash tc_1, \ldots, tc_k \cong tc'_1, \ldots, tc'_k}
	% 	\quad 
	% 	\inferrule*[right=cong. (f)]
	% 	{E \vdash tc \cong tc'}
	% 	{E \vdash f(tc) \cong f(tc')}
	% \end{mathpar}
\begin{mathpar}
    \inferrule*[right=start]
    {a = b \in E}
    {E \vdash a \cong b}
    \quad
    \inferrule*[right=cong. (variadic)]
    {E \vdash tc_i \cong tc'_i,\; i = 1, \ldots, k}
    {E \vdash tc_1, \ldots, tc_k \cong tc'_1, \ldots, tc'_k}
    \and
    \inferrule*[right=cong. (f)]
    {E \vdash tc \cong tc'}
    {E \vdash f(tc) \cong f(tc')}
\end{mathpar}
  Recall that we interpret the e-graph congruence by putting the equivalent terms together with a $\join$.
  The rule $tc \cong tc' \implies f(tc) \cong f(tc')$ is then exactly the property of composition in an $\catname{SLat}$-category, namely $f(tc \join tc') = f(tc) \join f(tc')$.
	Slotted e-graphs extend the above rules with the following:
	\begin{equation}%
		\label{eq:cong-bind}
		\inferrule*[right=cong. (bind)]
		{E \vdash tc \cong tc'}
		{E \vdash \text{bind}\; \$x\; tc \cong \text{bind}\; \$x\; tc'},
	\end{equation}
	plus the rules for equivalence of terms up to renaming of slots.
  This rule is exactly the property of preservation of joins by abstraction in~\eqref{eq:lambda-join-structural}.
  We did not have to introduce this rule in an ad hoc manner, nor did we have to justify it---it is provided by the axioms of the category we work in and thus the rule is sound with respect to semantics.

  The same restrictions apply as in the case of traditional e-graphs: if we were to draw a precise correspondence between string diagrams for $\catname{SLat}$-categories and slotted e-graphs, we would need to be careful about sharing and cycles.
  That is, we may introduce a comonoid generator, but it must not be natural on join morphisms.
  We argued that string diagrams may serve as a better and more principled representation in certain applications, but it is worth to note that slotted e-graphs are much more general.
  In particular, string diagrams require closed structure on the underlying category whereas slotted e-graphs may represent binding in theories that have no known closed structure, such as the $\pi$-calculus.



\section{Conclusion}
In this chapter we introduced string diagrams for $\catname{SLat}$-enriched symmetric monoidal closed categories and formalised them as concrete combinatorial objects.
The representation is shown to be sound and complete with respect to representation of closed-monoidal terms with join modulo $\equiv_{\mathsf{str}}$ which includes the SMC equations, $\catname{SLat}$-equations and the $\catname{SLat}$-rule for $\Lambda$-abstraction.
The equivalence of rewriting of closed-monoidal terms modulo $\equiv_{\mathsf{str}}$ and convex EDPOI rewriting was shown.
We also highlighted representational advantages of such string diagrams in the context of $\lambda$-calculi.
An important future work is exploring whether this combinatorial representation may lead to better algorithmics for doing equality saturation for $\lambda$-calculus.